\begin{problem}[Original selected-I3 auxiliary-array reconstruction]
\label{iv:open:K111:original_auxiliary_arrays}
\StatementRecord{iv:open:K111:original_auxiliary_arrays}{open}{}{the original common and endpoint arrays, ordered lower source chains, period maps, normal bundles and caps reproduced below; \cref{night:P3K.111} and its complete new auxiliary construction}
Reconstruct the originally reported 136,442-cone common atlas, its
three complete endpoint arrays and the ordered 324-chart lower source
chains, retaining all original period words, point/O modules, normal
derivatives and boundary caps.  The complete original files and their
ordered execution histories were not present in the supplied source.
The following preserved alternate construction therefore remains
unestablished at that exact numerical scope.  Its reported compact
pair $(2723,19124)$, lower total $(-2712,-19092)$ and nine-blowup
moving subtotal are not used in the proved construction of
\cref{night:P3K.111}.  That theorem uses the complete new atlas and
the actual source/exterior/cap maps of its supporting proof 20.

The required reconstruction includes the immutable zero upper
triangulation and the original ordered common-edge chain: the printed
support and band restrictions admit the explicitly exhibited diagonal
flip and do not determine those arrays.  Replaying them must recover
the whole fixed-owner partition and every original curve, surface,
point, normal and relative tangent contribution.  Each text below is
preserved with its original formula tags, at this stated unestablished
alternate-construction scope.
\end{problem}


\par\bigskip\noindent\textbf{Preserved alternate construction: original assembly introduction}\par
\noindent\textbf{Complete original-source attachment and directed physical transfer.}\par\smallskip

The nineteen complete supporting proofs are reproduced in their locked
order below. SI2 and I2P give the original source, the entire Miller and
theta coefficients, the point/O graphs and their inverses. I2M, I2Z,
I2O, I2I, I2F, I2ZC and I2L give the exact period dictionary, all three
native parents, the 136442-cone common fan and its proper endpoint lifts,
with full source bands, semigroup maps and coefficient inverses.
I2CV and CFL then calculate the complete compact fixed-stratum pair.

RB retains the whole scalar reference body and its three vanishing discs.
NS, NC and MN determine the two extending native charts, the actual cap
maps and their inverses, and the separate collision and reference normal
bundles. The class expressions on the closed replacement in NC.4 and MN
are transferred to the actual cap-completed bundles by the displayed
five-dimensional bundle bordism. IL retains all twelve lower/reference
doubles, including the intrinsic cube divisor and its boundary connection.
RFB and CSN give the full finite reference-fan continuation and independent
native scalar maps. ASI2 carries these maps through the compact attachment,
parameter square and original outer transfer.

The signed assembly is explicit. In the retained positive-inner
orientation, the compact pair is \((2723,19124)\), the twelve lower
pairs sum to \((-2712,-19092)\), and the difference of the two own-cap
moving pairs is \((-10,-28)\). Their sum is \((1,4)\), with the actual
RB reference in fan \(T_1\). RFB constructs the whole Jin reference
change \(J_{r,-1}-J_{r,1}=(-72,0)\), whose positive-inner counterpart
is \((72,0)\). Subtracting that reference gives \((-71,4)\); the
explicit Jin interval reversal gives \((71,-4)\). The original
physical-line dual and held-quaternion map then give
\(k_M=\mu+\nu\), \(\ell_M=4\mu\), hence \((-1,72)\).

The supporting sources retain their original progress and local-scope
sentences as part of their recorded proofs. The completed seed assertion
is ASI2. In particular RB's original point-coordinate descriptions are
carried to the two extending native charts by NS and MN; their finite
whole-pair comparison is the explicit RFB construction. These are the
maps used in the final assembly. The numbered dependencies are K68 for
the physical/Hopf inverse clutch, K96 for the original M2 good-root
attachment, F35 for the primitive common-link/Gysin/van Kampen maps,
and F41/F46 for the reference sphere construction. Their actual
hypotheses and boundary arrows are checked in the following proofs.


% Expanded verified cp36 input: notes/part3_e6_selected_i3_m2_return_support_01.tex


\par\bigskip\noindent\textbf{Preserved alternate construction: complete original auxiliary support 4}\par
\par\bigskip\noindent\textbf{P3K.111 supporting proof 4: Zero raw parent, lattice and full coefficients}\par\nobreak
\par

\medskip\noindent\textbf{Selected-I3 \(M_{2}\) source: the complete N12 zero endpoint}

\par\nobreak

This is the zero endpoint of SI2.1--6 for the actual line
\(M_2=\mathcal O([2G]-[O])\), translation by \(-3G\), original good
cap \((m,n)=(3,1)\), and coupled scalar
\(\lambda=s^{12}\lambda _0\), \(r_\lambda=s^4r_0\).
It constructs a finite regular endpoint with its two literal Tate source
bands, full period arrows, physical circle and original determinant.
The executed endpoint has 420 regular four-cones. The point/O modules
are carried on their actual graphs below. This note assigns no
characteristic pair or completed source return.

\par

\medskip\noindent\textbf{I2Z.1. Actual characters, units and the inspection lattice}

\par\nobreak

Use the SI2 characters and functions
\[
 X_0=V_{\rm reg}^{-1},\quad Y_0=Z,\quad p=R/P,\quad
 Q=P^3p,\quad
 f_0=(p-1)^{12}c_{12},\qquad
 c_{12}=\frac{\lambda _0c_q\eta^6}{u_s^{12}u_t^3}.
 \tag{I2Z.1}
\]
Here \(q_{\rm old}=Pc_qp\), \(s=P(p-1)/u_s\), \(t=Pu_t\), and
\[
 \eta=a_T\Theta'_Q(1)\frac{\Theta_Q(P^2)}{\Theta_Q(P)^2},
 \qquad \eta(0)=-1.
 \tag{I2Z.2}
\]
Thus the entire sixth power in (I2Z.1) includes
\(\Theta'_Q(1)^6\). All these functions are the original SI2 functions;
on a sufficiently small zero-end base disc \(p-1\), \(c_q,u_s,u_t,\eta\)
are nonzero. In particular \(f_0\) is an actual unit on this endpoint.
The original period words, without changing the physical line, are
\[
\begin{aligned}
 L_0=B^{-1}:&(X_0,Y_0)\longmapsto
       (P^{10}pf_0X_0,P^3Y_0),\\
 T_0=hB^{-1}:&(X_0,Y_0)\longmapsto
       (P^8pf_0X_0,P^6pY_0).
\end{aligned}
\tag{I2Z.3}
\]
The inverse arrows divide by the displayed factors at the same base
point. The base \(P,p\) is fixed by both periods, so these arrows commute,
including their units.

In the original valuation lattice
\(N=\mathbb Z^4_{(x,y,r,z)}
=\mathbb Z^4_{\operatorname{ord}(X_0,Y_0,P,p)}\), the period matrix is
\[
 M_0(r,z)=
 \begin{pmatrix}10r+z&8r+z\\3r&6r+z\end{pmatrix},
 \qquad \det M_0=36r^2+13rz+z^2.
 \tag{I2Z.4}
\]
The integer map
\[
 S=\begin{pmatrix}3&-3\\0&2\end{pmatrix},\qquad
 S\mathbb Z^2=3\mathbb Z\oplus2\mathbb Z
 \tag{I2Z.5}
\]
has index six. Its inverse is used only on that displayed image lattice.
It gives
\[
 SM_0=
 \begin{pmatrix}21r+3z&6r\\6r&12r+2z\end{pmatrix}.
 \tag{I2Z.6}
\]
Define \(\nu\) by affine interpolation of
\(\nu(n)=n(n-1)/2\) at every integer. Its slope on \((n,n+1)\)
is \(n\). The explicit convex potential
\[
 \Psi(\theta_1,\theta_2)=
 (15r+3z)\nu(\theta_1)+6r\nu(\theta_1+\theta_2)
       +(6r+2z)\nu(\theta_2)
 \tag{I2Z.6a}
\]
has the three stated positive slope jumps in the interior \(r,z>0\).
The identity
\(\nu(t+n)=\nu(t)+nt+n(n-1)/2\), valid on every interval
and at its endpoints, gives
\(\nabla\Psi(\theta+n)-\nabla\Psi(\theta)=SM_0n\).
The additive constant is exactly
\((15r+3z)n_1(n_1-1)/2+
6r(n_1+n_2)(n_1+n_2-1)/2+
(6r+2z)n_2(n_2-1)/2\).
At zero the subgradient is the sum of the three segments from zero
to \(-(15r+3z,0)\), \(-(6r,6r)\) and \(-(0,6r+2z)\).
All its vertices lie in the explicit lattice (I2Z.5).
Pulling this actual subgradient cell back by \(S^{-1}\), the
fundamental hexagon in the original lattice has vertices
\[
\begin{aligned}
 &(0,0),\ (-5r-z,0),\ (-10r-z,-3r),\\
 &(-13r-2z,-6r-z),\ (-8r-z,-6r-z),\
 (-3r-z,-3r-z).
\end{aligned}
\tag{I2Z.7}
\]
Direct substitution gives area \(\det M_0(r,z)\). Opposite edges
are paired by (I2Z.3), with the equations in I2Z.4 below. Thus the
index-six inspection map does not replace the original period lattice.

For comparison, the preserved exponent-six attempt had
\[
 M_r=\begin{pmatrix}4&2\\3&6\end{pmatrix},\quad
 M_z=\begin{pmatrix}1&1\\0&1\end{pmatrix},\quad
 M_z^{-1}M_r=\begin{pmatrix}1&-4\\3&6\end{pmatrix}.
 \tag{I2Z.8}
\]
Its characteristic polynomial is \(t^2-7t+18\), with discriminant
\(-23\). A constant real multiplier making both \(M_r,M_z\)
symmetric positive would make \(M_z^{-1}M_r\) self-adjoint for the
positive form supplied by \(M_z\), contradicting these nonreal roots.
This records the precise stopped constant-symmetrizer construction.
For the present exponent twelve the corresponding eigenvalues are
four and nine, and the actual map is (I2Z.5).

\par

\medskip\noindent\textbf{I2Z.2. The full raw endpoint and its invariant monomial maps}

\par\nobreak

At \(z=0,r=1\), retain the 72 triangles of the immutable upper file
\path{e6_selected_i3_m2_zero_upper.json}. Their vertices lie in the hexagon
\[
 (0,0),(-5,0),(-10,-3),(-13,-6),(-8,-6),(-3,-3),
 \tag{I2Z.9}
\]
and every triangle is primitive. Their union has area 36. This
triangulation includes the actual source line \(y+3r=0\).

Write the six lifted vertices as \(U_0,\ldots,U_5\), with
\((r,z)=(1,0)\), and set
\[
 A=(0,0,0,1),\ B_1=(-1,0,0,1),\
 B_2=(-2,-1,0,1),\ B_3=(-1,-1,0,1).
 \tag{I2Z.10}
\]
The following is the complete raw four-cone list, with each displayed
column set ordered to have positive determinant:

\begin{enumerate}
\def\labelenumi{\arabic{enumi}.}

\item
  \(A\) joined to each of the 72 upper triangles;
\item
  \((A,B_1,U_1,U_2)\);
\item
  \((A,B_1,(-10-k,-3-k,1,0),(-11-k,-4-k,1,0))\),
  \(k=0,1,2\);
\item
  \((A,B_1,U_3,B_2)\);
\item
  \((A,B_3,(-13+k,-6,1,0),(-12+k,-6,1,0))\),
  \(k=0,\ldots,4\);
\item
  \((A,B_3,U_3,B_2)\);
\item
  \((A,B_3,U_4,U_5)\).
\end{enumerate}

These 84 cones are obtained by coning the facets of the convex
\(r+z=1\) slice not containing \(A\). On each such facet the displayed
segments give its triangulation. Consequently the cones have disjoint
interiors and cover exactly the slice of (I2Z.7). The independent volume
identity is
\[
 \sum_\sigma\frac{|\det M_\sigma|}
 {6\prod_{\rho\subset\sigma}(r_\rho+z_\rho)}
 =\int_0^1\{36r^2+13r(1-r)+(1-r)^2\}\,dr
 =\frac{29}{2}.
 \tag{I2Z.11}
\]
Their determinants are one, two or three.

For a raw matrix \(M_\sigma\) of determinant \(d\), let \(u_1,\ldots,u_4\)
be its finite cover coordinates. The actual affine quotient functions
are the invariant monomials \(u^a\) with
\[
 a\in\mathbb Z_{\ge0}^4,\qquad aM_\sigma^{-1}\in\mathbb Z^4.
 \tag{I2Z.12}
\]
The saved list contains all such nonconstant monomials with
\(|a|\le d\). It generates the full invariant semigroup: an invariant
word of length greater than \(d\) in the character group of order \(d\)
has a nonempty zero-sum subword of length at most \(d\), by repetition
among its partial sums; remove it and continue. This supplies the
actual coordinate ring and finite quotient map.

For every refined cone with ray matrix \(N_\tau\), the recorded map is
\[
 u^a\longmapsto
 \prod_{\ell=1}^4 t_\ell^{(aM_\sigma^{-1}N_\tau)_\ell}.
 \tag{I2Z.13}
\]
Every exponent is an integer and nonnegative. On the dense torus its
inverse uses the same original four characters. This proves the
holomorphic morphism to the raw quotient, including its finite
isotropy, without adjoining fractional cover coordinates.

\par

\medskip\noindent\textbf{I2Z.3. Source-compatible resolution and the complete finite execution}

\par\nobreak

Insert, in this order, the three actual rays
\[
 (-9,-2,1,1),\quad(-7,-1,1,1),\quad(-6,-5,1,1).
 \tag{I2Z.14}
\]
The first two resolve the determinant-three raw cone; the third
resolves the determinant-two raw cone. The saved nonnegative
barycentric columns give each source and target monomial map through
(I2Z.13). All resulting cones have determinant one. All three inserted
rays have both base heights positive, so the upper 72 triangles remain
literally unchanged.

Now use the full source functional
\[
 \ell(x,y,r,z)=y+3r+z.
 \tag{I2Z.15}
\]
For every edge with opposite \(\ell\)-signs, the positive endpoint in
this execution has \(\ell=1\). Insert the sum of its two primitive
endpoints. In a regular chart this is the ordinary two-coordinate
blowup: on its two charts the original pair is \((a,ab)\) and
\((cd,d)\). Its inverse on the overlap divides by the displayed
nonzero coordinate. The negative endpoint value increases by one
along the resulting crossing edge, so the recorded finite sequence
ends with every cone on one side of (I2Z.15).

Complete the boundary period orbits of these rays in the already
subdivided fan. The order is material: an earlier attempt restarted
from the coarse regular fan and then inserted the same orbit members;
its executed cones still crossed (I2Z.15). That exact failed output is
preserved. The retained construction inserts into the actual
source-compatible fan, and every subsequent cone is contained in its
predecessor.

The resulting opposite period facets initially had different
diagonals. Their exact common subdivision is also retained. Normalize
each boundary ray by \(r+z\). On the \(L_0\) and \(L_0/T_0\) faces use
coordinates \((y/(r+z),r/(r+z))\); on the \(T_0\) faces use
\((x/(r+z),r/(r+z))\). Intersect all edges of a source triangulation
with all edges of the inverse-transported target triangulation using
the two-by-two rational determinant formula. Clear denominators and
divide by the greatest common divisor to obtain each actual primitive
intersection ray. Insert its full period orbit.

These insertions are performed through ordinary face blowups. In a
regular owner cone express the desired ray as an integral nonnegative
column. Insert the sum of the basis rays in its positive support,
then repeat in the child containing the desired ray. The sum of its
positive coefficients strictly decreases until that ray is present.
Every such insertion has the usual monomial blowup charts. The saved
execution has two face-overlay rounds: six new intersection rays
requiring twelve ordinary stars, then four requiring eight. The next
complete comparison has no mismatched period facet.

The full sequence has 77 inserted centres and 420 final regular
four-cones. Every final cone lies in its recorded raw parent and
still lies on one side of (I2Z.15). The exact parent-by-parent slice
volumes agree with their raw-parent volumes. The final incidence
table has 760 internal facets, 72 protected generic facets, two lower
facets and 86 period facets. These lists exhaust the boundary.

\par

\medskip\noindent\textbf{I2Z.4. Full boundary periods, source bands and point/O frames}

\par\nobreak

On valuation columns the complete two periods are
\[
 T_L=\begin{pmatrix}1&0&10&1\\0&1&3&0\\0&0&1&0\\0&0&0&1\end{pmatrix},
 \quad
 T_T=\begin{pmatrix}1&0&8&1\\0&1&6&1\\0&0&1&0\\0&0&0&1\end{pmatrix}.
 \tag{I2Z.16}
\]
Their three paired boundary types are
\[
\begin{array}{c|c|c}
\text{word}&\text{source equation}&\text{target equation}\\ \hline
L_0&x-y+7r+z=0&x-y=0\\
T_0&y+6r+z=0&y=0\\
L_0T_0^{-1}&3x-5y+15r+3z=0&3x-5y-6r-2z=0 .
\end{array}
\tag{I2Z.17}
\]
The full word \(L_0^aT_0^b\) adds the coefficient \(f_0^{a+b}\)
in the \(X_0\) character. If source and target ray matrices are
\(N_s,N_t\), the actual local coefficient arrow is
\[
 t'_i=f_0^{(N_t^{-1}(a+b,0,0,0)^t)_i}
       \prod_jt_j^{(N_t^{-1}T_L^aT_T^bN_s)_{ij}}.
 \tag{I2Z.18}
\]
The inverse uses \(N_s^{-1}T_T^{-b}T_L^{-a}N_t\) and
\(N_s^{-1}(-a-b,0,0,0)^t\). On the indicated localized faces all
exponents of vanishing coordinates are nonnegative integers. The
two matrices compose to identity, and the transported two unit
columns sum to zero. Thus every original \(p\), \(c_q\), \(u_s\),
\(u_t\), \(\lambda_0\) and \(\eta^6\) factor remains in the actual
forward and inverse maps.

For a final cone on \(\ell\ge0\), take \(j=1\); on \(\ell\le0\),
take \(j=2\). The literal source node coordinates are
\[
 x_j=Q^j Z,\qquad y_j=Q^{1-j}Z^{-1},\qquad x_jy_j=P^3p .
 \tag{I2Z.19}
\]
Their exponent rows on a ray \((x,y,r,z)\) are
\[
 y+3jr+jz,\qquad -y+3(1-j)r+(1-j)z.
 \tag{I2Z.20}
\]
They are nonnegative on every cone assigned to that band. Across
the band overlap \(x_2=Qx_1\), \(y_2=Q^{-1}y_1\), which is the actual
Tate node reindex. No full \(p\) power is suppressed.

Write \(f_b(x)=\Theta_Q(x/b)/\Theta_Q(x)\), \(b=P^2\).
The original \(M_{2}\) coefficient and the genuine point coefficient satisfy
\[
 V_{\rm reg}=\eta^{-1}P^{-2j}f_b(x_j)v_2,\qquad
 v_2=\frac{\eta P^{2j}}{X_0f_b(x_j)},\qquad
 V_{\rm pt}=\eta V_{\rm reg}.
 \tag{I2Z.21}
\]
For the actual point graph \(\operatorname{Bl}_{(x_j,P^2)}\), use
\(x_j=P^2a,\ ay_j=R\), or \(P^2=x_jc,\ y_j=Rc\).
The full SI2 theta product \(E\) gives the coefficient arrows
\[
 V_{\rm reg}=-\eta^{-1}P^{-2j}EA,\qquad
 V_{\rm reg}=\eta^{-1}P^{-2j}c^{-1}EC,\qquad C=-cA.
 \tag{I2Z.22}
\]
On the O chart \(v_2=(x_j-1)C_O\), the arrow is
\[
 V_{\rm reg}=\eta^{-1}P^{-2j-2}(x_j-P^2)E_O C_O .
 \tag{I2Z.23}
\]
The units \(E,E_O\), both graph charts, the section coincidence and
their inverse coefficient maps are those proved in SI2.3. Equations
(I2Z.19)--(I2Z.23) are the source and line morphisms carried by this
endpoint. Finite marked-boundary graph pullbacks use precisely these
ideals and module generators.

\par

\medskip\noindent\textbf{I2Z.5. Physical weights, original determinant and certificate scope}

\par\nobreak

The original physical circle has weight
\((-1,0,0,0)\) in the \((X_0,Y_0,P,p)\) characters. Therefore in
every final chart \(N_\tau\) its exact coordinate weights are
\[
 N_\tau^{-1}(-1,0,0,0)^t.
 \tag{I2Z.24}
\]
The nonzero weights, including every finite stabilizer they determine,
are retained in the output. Period arrows (I2Z.18) commute with this
circle because their coefficients are source-base functions.

With the original SI2 differential and
\(J_{ts}=\det\partial(t,s)/\partial(P,R)\), the original four-form is
\[
 \Omega_{\rm original}
 =\frac{2J_{ts}}{a_T}\,
 P\,dP\wedge dp\wedge d\log X_0\wedge d\log Y_0,\qquad
 J_{ts}(0)=\frac12,\quad a_T(0)=1 .
 \tag{I2Z.25}
\]
The character order and the wedge order in (I2Z.25) have positive
comparison. Each saved ray matrix has determinant \(+1\); hence a
ray \((x,y,r,z)\) has original canonical order \(2r+z-1\).
The unit \(2J_{ts}/a_T\), including its phase, is retained. This is
the original determinant map through the base blowup and the actual
source reindices.

The executable \path{e6_selected_i3_m2_zero_fan.py} proves 4,320 exact
identities and writes the full cone, invariant-semigroup, source,
point-module, unit, inverse, physical-weight and canonical-order
tables. Its input upper file is immutable. Its proof/source hashes
are locked on execution. Every stellar centre and every period-face
intersection is saved, so the finite construction is reproducible.
The two stopped source/period attempts are preserved under
\path{records/attempts/e6_selected_i3_m2_zero_first_source_orbit} and
\path{records/attempts/e6_selected_i3_m2_zero_before_period_face_overlay.}
The resulting endpoint is the zero parent for the subsequent common
generic fan lift; its global compact fixed-locus and outer transfer
are computed only after those full three-end identifications.


% End expanded verified cp36 input


% Expanded verified cp36 input: notes/part3_e6_selected_i3_m2_return_support_05.tex


\par\bigskip\noindent\textbf{Preserved alternate construction: complete original auxiliary support 5}\par
\par\bigskip\noindent\textbf{P3K.111 supporting proof 5: One raw parent and full coefficients}\par\nobreak
\par

\medskip\noindent\textbf{Selected-I3 \(M_{2}\): the complete N12 one parent}

\par\nobreak

We retain the original source \(y^{2}=x^{3}+r_{\rm source} x^{2}+t x+t^{2}\), \(G=(0,t)\),
\(M_{2}=\mathcal O([2G]-[O])\), translation \(-3G\), \(s=1/4-r_{\rm source}\) and the selected
height \(q=t-t_c(s)\), \(t_c(0)=0\). SI2 fixes the actual coupled scalar
\(\lambda=\lambda_0 s^{12}\), its root \(r_{\rm sc}=r_{0} s^{4}\), and the full source,
finite character and four-root-turn reference square. This note
constructs the one parent in that original lattice.

\par

\medskip\noindent\textbf{I2O.1. Complete period functions and the original lattice}

\par\nobreak

Use the entire SI2 functions \(Q=P^{2}R\), \(R-P=s u_s\), \(t=P u_t\) and
\(q=R c_q\). Put \(p=R/P\), \(v=p-1\) and

\[
\begin{gathered}
c_{12}=\lambda_0 c_q \eta^{6}/(u_s^{12} u_t^{3}), \\
f_{0}=v^{12} c_{12}.
\end{gathered}
\]

The genuine-coordinate periods in order \((V_{\rm reg},Z)\) are

\[
\begin{gathered}
K=hB=(P^{12} p v^{12} c_{12} V_{\rm reg},p Z), \\
J=h=(P^{2} V_{\rm reg},P^{3}p Z).
\end{gathered}
\]

These retain their original point-line module \(V_{\rm reg}=\eta^{-1} V_{{\rm pt},2}\)
and the full theta identity with \(v_{2}\). Set \(X=Z^{-1},Y=V_{\rm reg}/Z\) and
\(L=J^{-1},T=K\). Their exact arrows are

\[
\begin{gathered}
L:(X,Y)\longmapsto (P^{3}p X,Pp Y), \\
T:(X,Y)\longmapsto (p^{-1} X,P^{12} v^{12} c_{12} Y).
\end{gathered}
\]

The character inverse is \(Z=X^{-1},V_{\rm reg}=Y/X\); the period inverse
is \(J=L^{-1},K=T\). Both matrices are integral and determinant 1.
Every inverse coefficient is the reciprocal of the displayed
full function. These functions are invariant under all periods.
For \(r=\operatorname{ord} P,z=\operatorname{ord} v\) the exact period matrix and determinant are

\[
\begin{gathered}
N=\begin{pmatrix}3r&0\\r&12r+12z\end{pmatrix}, \operatorname{det} N=36r^{2}+36rz.
\end{gathered}
\]

The inspection map \((x,y)\longmapsto (x,-x+3y)\) maps the original integer
lattice isomorphically onto its explicit index 3 image. It sends
N to \(\operatorname{diag}(3r,36r+36z)\). The rectangular quadratic cells for the
form \(3r \binom{n_{1}}{2}+(36r+36z)\binom{n_{2}}{2}\), pulled back through
this specified lattice map, have the original parallelogram

\[
\begin{gathered}
(0,0),(3r,r),(3r,13r+12z),(0,12r+12z).
\end{gathered}
\]

Its original area equals \(\operatorname{det} N\). Its opposite faces are paired
by exactly L and T, including the full p and \(c_{12}\) coefficients.
The finite-index image is retained throughout this construction.

\par

\medskip\noindent\textbf{I2O.2. Integral parent charts and proper regular subdivision}

\par\nobreak

At \(r=1,z=0\) triangulate the upper parallelogram by inserting its
integral points into every containing triangle or edge, starting
with the two diagonal triangles. The executed upper list has 72
primitive triangles, area 36, with matching period boundary edges.
At \(r=0\) use the thirteen primitive lower rays \(E_j=(0,j,0,1)\),
\(0\le j\le 12\). Join \(E_0\) to all upper triangles. For \(0\le j<12\) join
\(E_j\),\(E_{j+1}\) to the upper edge from \((3,13,1,0)\) to \((0,12,1,0)\).
These lower parents have determinant 3. The joins cover every
slice of the displayed parallelogram, and their common faces
are exactly the shared faces of these joins.

For each singular primitive matrix B, the code enumerates the
finite residue group \(B^{-1} \mathbb Z^4/\mathbb Z^4\) using its four inverse columns.
Its cardinality equals \(\operatorname{det} B\). Each nonzero representative gives
an integral point of the half-open fundamental parallelepiped.
Divide it by its gcd and insert the selected primitive ray in
every cone containing it. The positive barycentric coefficients
are strictly less than 1, so every changed determinant strictly
decreases. Every step retains its full parent and child list.
The process terminates in integral determinant-one cones.

For each raw parent the primitive dual rays and all lattice
points of their half-open parallelepiped generate its entire
dual semigroup: subtract the coefficient floors of an arbitrary
point. The certificate records this entire finite generating
list and all nonnegative integral exponent maps to the regular
charts. On their original dense torus the inverse is the exact
rational character matrix. The subdivisions preserve the whole
support; these affine maps therefore assemble to the proper
map from the regular parent to the original toric support.

\par

\medskip\noindent\textbf{I2O.3. Full boundary arrows, source modules and volume}

\par\nobreak

The complete periodic facets are

\[
\begin{gathered}
L: x=0 \longmapsto  x=3r, \\
T: 3y-x=0 \longmapsto  3y-x-36r-36z=0.
\end{gathered}
\]

Their original lattice shifts are \((3r,r)\) and \((0,12r+12z)\).
The full p columns are \((1,1)\) and \((-1,0)\), respectively; the \(c_{12}\)
columns are \((0,0)\) and \((0,1)\). The inverse maps use the opposite
shifts and reciprocal units. The executed full sets of final
three-faces agree under these exact maps. For each regular
source/target pair their matrix is \(B_{\rm target}^{-1} T_{\rm period} B_{\rm source}\)
and each full unit column is \(B_{\rm target}^{-1} u_{\rm period}\). Composing
with the actual inverse cancels both the matrix and each full
coefficient, giving the original quotient gluing.

Every source-band chart has the full functions

\[
\begin{gathered}
x_{\rm source}=QZ=pP^{3}/X, y_{\rm source}=Z^{-1}=X, \\
x_{\rm source} y_{\rm source}=Q.
\end{gathered}
\]

Their ray orders \(3r-x\) and x are nonnegative on every cone.
The full inverse Tate chart therefore extends to the original
Weierstrass source. The actual original point-line coefficient
and its inverse arrow are

\[
\begin{gathered}
v_{2}=\eta P^{2} Y/[X f_{P^{2}}(pP^{3}/X)], \\
Y=\eta^{-1} P^{-2} X f_{P^{2}}(pP^{3}/X) v_{2},
\end{gathered}
\]

as morphisms in the specified point/O Cartier modules. The
proper graph of \((x_{\rm source},P^{2})\) is carried through these exact
arrows, with its point and zero-section factors retained.

The physical circle has valuation vector \((0,1,0,0)\); its four
weights on each regular chart are the inverse matrix applied
to this vector. In the original source-volume coefficient the
change to \((X,Y)\) gives exactly

\[
\begin{gathered}
\Omega=-(2J_{\rm ts}/a_T) P \mathrm dP\wedge \mathrm dv\wedge \mathrm d\log X\wedge \mathrm d\log Y.
\end{gathered}
\]

The Jacobian \(J_{\rm ts}\) and invariant differential \(a_T\) remain their
entire nonzero units. Every original canonical ray order is
\(2r+z-1\). The coupled scalar contributes only source-base forms
to \(\mathrm d\log V\), whose wedge with \(\mathrm dP\wedge \mathrm dv\) vanishes. All source orders,
physical weights, inverse maps and canonical orders are saved
in the certificate. The characteristic pair awaits the common
attachment, complete point graph and compact normal ledger.

The \(\eta\) in every displayed point-module arrow is exactly
\(a_T \Theta'_Q(1) \Theta_Q(P^{2})/\Theta_Q(P)^{2}\), with \(\eta(0)=-1\)
and \(a_T(0)=1\). In \(c_{12}\) its full sixth power includes
\(\Theta'_Q(1)^{6}\). The source SI2 certificate retains this entire
identity and the reciprocal frame morphism.


% End expanded verified cp36 input


% Expanded verified cp36 input: notes/part3_e6_selected_i3_m2_return_support_06.tex


\par\bigskip\noindent\textbf{Preserved alternate construction: complete original auxiliary support 6}\par
\par\bigskip\noindent\textbf{P3K.111 supporting proof 6: Infinity raw parent and full coefficients}\par\nobreak
\par

\medskip\noindent\textbf{Selected-I3 \(M_{2}\): the complete N12 infinity parent}

\par\nobreak

We retain the original source \(y^{2}=x^{3}+r_{\rm source} x^{2}+t x+t^{2}\), \(G=(0,t)\),
\(M_{2}=\mathcal O([2G]-[O])\), translation \(-3G\), \(s=1/4-r_{\rm source}\) and the selected
height \(q=t-t_c(s)\), \(t_c(0)=0\). The actual balanced and genuine point coefficients are
\(V_{2,T}=Z^{-1} f_P(Z)^{2} r_{2} v_{2}\) and
\(V_{\rm reg}=Z V_{2,T}=\eta^{-1} V_{{\rm pt},2}\), with the full OG/MP point/O module
map retained. For this actual coupled
filling choose \(\lambda=\lambda_0 s^{12}\) and root \(r_{\rm sc}=s^{4} r_{0}\),
\(r_{0}^{3}=\lambda_0\). The full four-root-turn reference square is carried
by the source proof. This note constructs its infinity source,
period and toric parent, without assigning a characteristic pair.

\par

\medskip\noindent\textbf{I2I.1. Entire functions and the original lattice}

\par\nobreak

Use \(Q=P^{2}R\), \(R-P=s u_s\), \(t=P u_t\), \(q=R c_q\), with the original
analytic units of E81/SI3. Put

\[
\begin{gathered}
c_{12}=\lambda_0 c_q \eta^{6}/(u_s^{12} u_t^{3}), \\
p=R/P, f_{0}=(p-1)^{12} c_{12}.
\end{gathered}
\]

The full genuine-coordinate periods, in order \((V_{\rm reg},Z)\), are

\[
\begin{gathered}
K=hB=(P^{12} p f_{0} V_{\rm reg},pZ), \\
J=h=(P^{2} V_{\rm reg},P^{3}pZ).
\end{gathered}
\]

They fix P,p and every source function. At infinity set \(P=Ru\),
\(p=u^{-1}\) and \(f=(1-u)^{12} c_{12}(P=Ru,R)\). The exact unit equation is
\(f_{0}=u^{-12} f\). Thus the actual periods are

\[
\begin{gathered}
K=(R^{12} u^{-1} f V_{\rm reg},u^{-1} Z), \\
J=(R^{2}u^{2} V_{\rm reg},R^{3}u^{2} Z).
\end{gathered}
\]

Every inverse uses these reciprocal full functions at the actual
inverse point. With \(X=V_{\rm reg}/Z\) and \(L=K^{2}J,T=KJ\), direct composition
gives

\[
\begin{gathered}
L:(X,Z)\longmapsto (R^{23} f^{2} X,R^{3} Z), \\
T:(X,Z)\longmapsto (R^{11} f X,R^{3}u Z).
\end{gathered}
\]

The inverse period maps are \(K=L T^{-1}\) and \(J=T^{2} L^{-1}\); the character
inverse is \(V_{\rm reg}=XZ\). The two integer matrices have determinant 1.
The valuation matrix for \(r=\operatorname{ord} R,z=\operatorname{ord} u\) is

\[
\begin{gathered}
N(r,z)=\begin{pmatrix}23r&11r\\3r&3r+z\end{pmatrix}, \\
\operatorname{det} N=36r^{2}+23rz.
\end{gathered}
\]

The inspection map \((X_{\rm val},Z_{\rm val})\longmapsto (3X_{\rm val},11Z_{\rm val})\) is an isomorphism
from the original integer lattice onto its stated index 33 image.
On that image the period matrix is the symmetric matrix
\(\begin{pmatrix}69r&33r\\33r&33r+11z\end{pmatrix}\). Its actual negative-square quadratic
on period labels is
\(69r \binom{n_{1}}{2}+(33r+11z)\binom{n_{2}}{2}+33r n_{1} n_{2}\).
Its affine interpolation has transverse jumps \(36r,33r,11z\).
The full change under a label shift is the bilinear form of
that matrix. Pulling its dual cells back by the specified
lattice isomorphism gives the original-lattice hexagon

\[
\begin{gathered}
(0,0), (-12r,0), (-23r,-3r), \\
(-23r,-3r-z), (-11r,-3r-z), (0,-z).
\end{gathered}
\]

Its opposite faces are glued by \(L,T,L/T\). Its area is exactly
\(\operatorname{det} N\). At \(z=0\) the generic parallelogram has area 36. This
constructs the complete periodic support in the original lattice.

\par

\medskip\noindent\textbf{I2I.2. Full finite parent and actual regular subdivision}

\par\nobreak

Triangulate the upper parallelogram at \(r=1,z=0\) by successively
inserting all its integral points into every containing triangle
and edge, starting with its two diagonal triangles. The exact
output has 72 determinant-one triangles and total area 36. The
boundary subdivisions agree under both original period sides.
At \(r=0\) the lower segment has endpoints \(A_{0}=(0,0,0,1)\) and
\(A_{1}=(0,-1,0,1)\). Join \(A_{0}\) to all upper triangles. Join \(A_{0},A_{1}\) to
each of the twelve consecutive horizontal unit edges from
\((-23,-3)\) to \((-11,-3)\), and to the remaining diagonal edge from
\((-11,-3)\) to \((0,0)\). These cones cover the hexagonal support on
each base slice \(r,z\ge 0\) and meet on their stated common faces.
The horizontal joins have determinant 1 and the diagonal join
has determinant 11.

For each singular cone with ordered primitive matrix B, enumerate
\(B^{-1} \mathbb Z^4/\mathbb Z^4\) by the residues of its four inverse-coordinate
columns, modulo its determinant. The resulting finite group has
exactly \(\operatorname{det} B\) elements. For each nonzero residue its vector in
the half-open fundamental parallelepiped is integral. Divide by
its gcd, choose a primitive ray, and subdivide every cone
containing it. A positive barycentric coefficient is strictly
less than 1, so every changed child\textquotesingle s determinant is a positive
integer smaller than its parent\textquotesingle s. The algorithm terminates.
Every step is the actual toric proper subdivision, with its
full integer character maps, and affects the complete incident
star. The certificate retains all selected rays and children.

For an original raw cone enumerate the primitive dual rays and
every integral vector in their half-open parallelepiped. They
generate the whole dual semigroup: subtract the nonnegative
coefficient floors of an arbitrary dual-cone lattice point.
The recorded monomials for every final regular cone are the
pairings of those generators with its primitive columns. They
are integral and nonnegative. Thus each chart has its full
holomorphic morphism to the original affine semigroup chart,
including the determinant-eleven raw parent. On the dense
original torus the inverse is the stated rational integer
character map. Every subdivision preserves the whole support,
so their assembled map is proper.

\par

\medskip\noindent\textbf{I2I.3. All original period faces, point modules and volume}

\par\nobreak

On the final four-dimensional fan the full face arrows are

\[
\begin{gathered}
L: x+23r=0 \longmapsto  x=0, \\
T: y+3r+z=0 \longmapsto  y=0, \\
L/T: 3x-11y+36r=0 \longmapsto  3x-11y-11z=0.
\end{gathered}
\]

Their lattice shifts are respectively \((23r,3r),(11r,3r+z)\),
\((12r,-z)\); their full unit columns for f are \((2,0),(1,0),(1,0)\).
The inverse arrows retain the opposite shifts and reciprocal
units. The executed boundary face sets are exactly paired by
these maps. At each regular chart their character matrix is
\(B_{\rm target}^{-1} T_{\rm period} B_{\rm source}\), with full unit column
\(B_{\rm target}^{-1} u_{\rm period}\). Its inverse has the opposite transported
column; the coefficient product cancels identically. This
specifies the full original toric quotient gluing.

In every chart the actual source-band functions are

\[
\begin{gathered}
x_{\rm source}=QZ=R^{3}u^{2} Z, y_{\rm source}=Z^{-1}, \\
x_{\rm source} y_{\rm source}=Q.
\end{gathered}
\]

Their ray orders are \(y+3r+2z\) and \(-y\). They are nonnegative on
all final cones; their sum is the full source order \(3r+2z\).
The full inverse Tate map therefore carries these source-node
charts to the original Weierstrass source. On this band the
actual original point coefficient is

\[
\begin{gathered}
v_{2}=\eta P^{2} XZ/f_{P^{2}}(QZ), P=Ru,
\end{gathered}
\]

with its reciprocal arrow in the stated Cartier point/O modules.
This is the full MP theta identity, including every \(\eta\), P and
source-band factor. Its subsequent point graph is the proper
graph of \((x_{\rm source},P^{2})\), to be attached through these very maps.

The physical circle has valuation vector \((1,0,0,0)\); each chart\textquotesingle s
four physical weights are its full inverse matrix applied to
that vector. The source Jacobian and invariant differential
remain their complete units. In the common coordinate volume,
\(\Omega=-(2J_{\rm ts}/a_T) P \mathrm dP\wedge \mathrm dp\wedge \mathrm d\log X\wedge \mathrm d\log Z\), substitute
\(P=Ru,p=u^{-1}\). Since \(P \mathrm dP\wedge \mathrm dp=-R \mathrm dR\wedge \mathrm du\), the full volume is

\[
\begin{gathered}
\Omega=(2J_{\rm ts}/a_T) R \mathrm dR\wedge \mathrm du\wedge \mathrm d\log X\wedge \mathrm d\log Z.
\end{gathered}
\]

Every ray therefore has original canonical order \(2r+z-1\), with
the actual nonzero coefficient retained. The change from the
coupled scalar\textquotesingle s line factor adds only source-base differentials
to \(\mathrm d\log V\), which wedge to 0 with the two displayed source forms.
The certificate checks every source exponent, inverse matrix,
physical weight and canonical order. No global characteristic
pair is assigned by this parent construction.

The executed certificate \path{code/part3_root/e6_selected_i3_m2_infinity_fan.py}
passes 598 exact checks. Its raw upper slice has 72 regular triangles;
its complete regular infinity parent has 149 four-cones after 18
recorded primitive stellar insertions. Every original semigroup
generator, full inverse monomial map, physical weight, source order,
canonical order and period-face set is retained in its JSON.

The \(\eta\) in every displayed point-module arrow is exactly
\(a_T \Theta'_Q(1) \Theta_Q(P^{2})/\Theta_Q(P)^{2}\), with \(\eta(0)=-1\)
and \(a_T(0)=1\). In \(c_{12}\) its full sixth power includes
\(\Theta'_Q(1)^{6}\). The source SI2 certificate retains this entire
identity and the reciprocal frame morphism.


% End expanded verified cp36 input


% Expanded verified cp36 input: notes/part3_e6_selected_i3_m2_return_support_07.tex


\par\bigskip\noindent\textbf{Preserved alternate construction: complete original auxiliary support 7}\par
\par\bigskip\noindent\textbf{P3K.111 supporting proof 7: Complete generic common fan}\par\nobreak
\par

\medskip\noindent\textbf{Selected-I3 \(M_{2}\) N12: the complete generic common fan}

\par\nobreak

We retain SI2.1-\/-6 and the full period dictionary I2M.1-\/-3.
The original point line is \(M_{2}=\mathcal O(2G-O)\), and its source
translation is minus three G. This construction gives the
common generic boundary used by the three endpoint families.
The point-section graph extensions and their compact normal
ledger remain subsequent constructions.

\par

\medskip\noindent\textbf{I2F.1. The actual torus and the three parent morphisms}

\par\nobreak

Use the characters \((X_c,Z_c,P)=(V_{\rm reg}/Z,Z,P)\) and retain
the entire functions p and \(f_{0}\) on the punctured exceptional
parameter curve. The complete period maps are

\[
 L_c:(X_c,Z_c,P)\mapsto
 (P^{23}p^{-1}f_0^2X_c,P^3p^3Z_c,P),
\]
\[
 T_c:(X_c,Z_c,P)\mapsto
 (P^{11}p^{-1}f_0X_c,P^3p^2Z_c,P).
 \tag{I2F.1}
\]

They fix p and every coefficient in \(f_{0}\), so the maps commute
with their full analytic functions. The valuation lattice
has ordered basis \((\operatorname{ord} X_c,\operatorname{ord} Z_c,\operatorname{ord} P)\), and a period word
\(L_c^i T_c^j\) adds \(((23i+11j)h,(3i+3j)h,0)\) to a ray \((a,b,h)\).
Its full unit columns are

\[
 u_p(i,j)=(-i-j,3i+2j,0)^t,\qquad
 u_f(i,j)=(2i+j,0,0)^t.
 \tag{I2F.2}
\]

The original zero-parent character map is
\((X_c,Z_c)=(1/(X_{0}Y_{0}),Y_{0})\), with inverse
\((X_{0},Y_{0})=(1/(X_cZ_c),Z_c)\). The one-parent map is
\((X_c,Z_c)=(Y_{1},1/X_{1})\), with inverse
\((X_{1},Y_{1})=(1/Z_c,X_c)\). Their complete period changes
are the integral matrices I2M.2 and I2M.2, with the
written inverse matrices. At infinity the characters
already agree, while the transverse parameter is
the full function \(R=pP\). There \(f_{\rm infinity}=p^{-12} f_{0}\).
These identities give the full parent morphisms in
both directions on the common torus.

The generic period matrix is
\[
                   M=\begin{pmatrix}23&11\\3&3\end{pmatrix}.
 \tag{I2F.3}
\]
The polygon H with ordered vertices
(0,0),(-12,0),(-23,-3),(-11,-3) has area thirty-six.
I2I.1 proves that its translates by \(M \mathbb Z^{2}\) have
disjoint interiors and cover the plane. the zero-parent construction
and I2O.1-\/-3 give the other two fundamental domains
with the literal same period lattice through the
maps just specified. Each parent has seventy-two regular
upper triangles in its own lattice. Every one of
those triangles and its ordered rays remains in
the common-fan certificate.

\par

\medskip\noindent\textbf{I2F.2. Exact finite overlay and boundary closure}

\par\nobreak

For each zero-parent triangle, apply its character
matrix and every period translate that meets H.
The possible translate indices are bounded exactly
by the two coordinates \(M^{-1}(x,y)\): on each convex
polygon, their extrema occur at vertices. Taking
integer ceilings and floors of the differences of
the extremal intervals gives the complete finite
list. The code uses rational arithmetic for all
these bounds and intersections.

Intersect those translated triangles with the
seventy-two infinity-parent triangles in H. Then
intersect each nonempty polygon with every
translated one-parent triangle meeting it. Each
intersection is constructed by clipping against
the three oriented rational half-planes of the
triangle. The interpolation parameter is exactly
the ratio of the signed values of the defining
linear form at the two endpoints. This gives
6387 polygons with positive area; their areas sum
to thirty-six. Their interiors are disjoint because
each of the three original periodic triangulations
has that property.

Every polygon vertex is inserted in all incident
polygon edges. Its period images lying in H are
also inserted. The resulting vertex set has 6043
points. The two boundary identifications are
\[
 (a,b,h)\mapsto(a+11h,b+3h,h)
 \quad\hbox{from }b+3h=0\hbox{ to }b=0,
\]
\[
 (a,b,h)\mapsto(a+12h,b,h)
 \quad\hbox{from }3a-11b+36h=0
                  \hbox{ to }3a-11b=0.
 \tag{I2F.4}
\]
These are \(T_c\) and \(L_c T_c^{-1}\), with their full
coefficients I2F.1-\/-2. The certificate verifies
equality of the entire subdivided edge sets in
both identifications. These two words generate
the original period lattice.

Each convex intersection polygon is triangulated
by its ordered ears. A chosen ear has positive
oriented area and contains no other remaining
polygon vertex. Removing it preserves a polygon
with the same ordered boundary. The finite
recorded removals therefore partition the polygon
without overlapping interiors or an omitted edge.
All triangles are cones over primitive integer
rays: clear the two rational denominators and
divide by the gcd of the resulting triple.

\par

\medskip\noindent\textbf{I2F.3. The executed periodic regularization}

\par\nobreak

For a simplicial integer cone B of determinant d,
the coefficients of each ambient coordinate
vector in B give generators of the quotient
lattice \(B^{-1} \mathbb Z^{3}/\mathbb Z^{3}\). The code enumerates that
finite subgroup by exact addition modulo d,
and checks that it has precisely d elements.
Each nonzero residue gives an integral vector
in the half-open fundamental parallelepiped.
Its primitive ray and all its period images
in H are inserted in every containing cone.

The record retains every selected ray, all its
period images, the number of changed cones and
the selected determinant. On every changed cone
the new ray has nonnegative barycentric entries;
every positive entry is strictly less than one.
Replacing each positive-coordinate ray by the
inserted ray gives the ordinary stellar
subdivision. Its children cover the parent,
have disjoint interiors and agree on every
unchanged face. Their determinants are positive
integers strictly smaller than the parent\textquotesingle s.
The same operation on every incident cone and
period image retains the entire common boundary.

The exact regularization has completed: 62,281 recorded stellar
steps give 136,442 regular cones. Every final determinant is one.
There are 204,429 interior edges, each incident to exactly two cones,
and 468 boundary edges, each incident to exactly one cone before
period gluing. The exact height-one area is thirty-six. Both entire
period boundary edge sets agree under I2F.4. Every insertion history,
primitive height and original parent is retained in the certificate.
These computations establish the stated common generic atlas; no
numerical characteristic pair is assigned at this stage.

\par

\medskip\noindent\textbf{I2F.4. Every affine parent map with its complete units}

\par\nobreak

Let B be an ordered final ray matrix. Its
chart coordinates theta give
\[
                      (X_c,Z_c,P)=\theta^B.
 \tag{I2F.5}
\]
For each zero or one parent, the certificate
retains its translated ray matrix Bk, original
triangle index and exact word(i,j). The actual
parent coordinates are
\[
 \theta_{k,\alpha}=
 p^{-(B_k^{-1}u_p)_\alpha}
 f_0^{-(B_k^{-1}u_f)_\alpha}
 \prod_\beta\theta_\beta^{(B_k^{-1}B)_{\alpha\beta}}.
 \tag{I2F.6}
\]
Every entry of \(B_{k}^{-1} B\) is an integer at least
zero. All unit exponents are integers. The
affine formula therefore extends over the
complete child chart, and its inverse on the
original torus is the full character formula
I2F.1-\/-2 with the opposite word. The p and \(f_{0}\)
factors have their entire source-parameter
dependence in I2F.6.

For the infinity parent, let Bi denote its
original ordered ray matrix in (\(X_c\),\(Z_c\),R).
The literal parameter equality \(R=pP\) gives
\[
 \theta_{\infty,\alpha}=
 p^{(B_i^{-1}(0,0,1)^t)_\alpha}
 \prod_\beta\theta_\beta^{(B_i^{-1}B)_{\alpha\beta}}.
 \tag{I2F.7}
\]
Again all character exponents are nonnegative
integers, and the complete coefficient is a
unit on the generic exceptional domain. The
saved orientation permutation relates Bi to
the exact parent triangle order. These three
families of maps preserve every original
character. Character localizations give their
overlap morphisms, with the same full functions
and inverse period words. The subdivisions
cover each parent support; hence these affine
maps assemble to the proper toric refinements
of the three specified generic parent atlases.

\par

\medskip\noindent\textbf{I2F.5. Original source, line, physical circle and volume}

\par\nobreak

The original source node coordinates are
\[
 Z_1=pP^3Z_c,\qquad W_1=Z_c^{-1},\qquad
 Z_1W_1=Q=pP^3.
 \tag{I2F.8}
\]
On a ray \((a,b,h)\), their exponents are b+3h
and minus b. Both are nonnegative everywhere
on H. The coefficient p in the first coordinate
is retained. E81\textquotesingle s full inverse Tate map carries
these node coordinates to the original
Weierstrass source. The point-line morphism is
\[
 v_2=\frac{\eta P^2X_cZ_c}{f_{P^2}(pP^3Z_c)}.
 \tag{I2F.9}
\]
This follows from the full SI2 point identity and the full theta
identity \(f_{P^{2}}(QZ)=P^{2} f_{P^{2}}(Z)\). Both the
point and zero Cartier modules, including the
reciprocal line arrow, remain specified in
I2F.9 for their subsequent common graphs.

The point coefficient here is the actual Miller unit
\(\eta=a_T \Theta'_Q(1) \Theta_Q(P^{2})/\Theta_Q(P)^{2}\). Its value at the
source origin is −1. Both forward and reciprocal maps retain it.

The physical circle cocharacter is (1,0,0),
so its exact three chart weights are \(B^{-1}\)
times that vector. The program checks every weight is integral and retains every
primitive ray height. No finite-index
change of the physical circle has been made.

The full original source volume becomes
\[
 \Omega=-\frac{2J_{ts}}{a_T}\,
 P\,dP\wedge dp\wedge d\log X_c\wedge d\log Z_c.
 \tag{I2F.10}
\]
The minus sign is the determinant of
(\(X_{0}\),\(Y_{0}\)) to (\(X_c\),\(Z_c\)), namely minus one.
The complete coefficient \(2J_{\rm ts}/a_T\) is the
nonvanishing analytic source coefficient retained in SI2 and I2I.3. Inserting I2F.5 gives original
canonical order 2h-1 along a ray of height h;
each such integer is stored in the certificate.

The executing program is
\path{code/part3_root/e6_selected_i3_m2_common_fan.py}. It retains each
ordered matrix, the three parent morphisms, every full unit vector,
the periodic stellar history and the locked source files. Its completed
certificate is \path{code/part3_root/e6_selected_i3_m2_common_fan.json},
SHA256 \path{2999175af858e837e202908b15f4743fd7b5caf40af321673a05925ebbfe98cf}.
Every support, incidence, integral inverse and original source-map
gate has passed. All five original source locks were independently
recomputed after the certificate closed and matched. The full
corner lifts, global point graphs and compact normal attachments
remain the next calculations.


% End expanded verified cp36 input


% Expanded verified cp36 input: notes/part3_e6_selected_i3_m2_return_support_08.tex


\par\bigskip\noindent\textbf{Preserved alternate construction: complete original auxiliary support 8}\par
\par\bigskip\noindent\textbf{P3K.111 supporting proof 8: Complete zero lift}\par\nobreak
\par

\medskip\noindent\textbf{Selected-I3 \(M_{2}\) N12: full zero-end lift of the common atlas}

\par\nobreak

This working proof retains the original source, selected height,
\(M_2=\mathcal O([2G]-[O])\), translation by \(-3G\), and
the cubic good-cap and four-root-turn maps of SI2. The source-compatible
zero parent is I2Z\textquotesingle s executed 420-chart fan. The common input is
\path{code/part3_root/e6_selected_i3_m2_common_fan.json}, with SHA256
\path{2999175af858e837e202908b15f4743fd7b5caf40af321673a05925ebbfe98cf}:
it contains 136,442 regular generic cones and all full parent labels.
The companion of this note computes the complete four-dimensional
lift and verifies its source, point, period and determinant maps.
No characteristic pair or return is assigned by this construction.

\par

\medskip\noindent\textbf{I2ZC.1. Actual characters and their full coefficients}

\par\nobreak

Use the original zero-end torus
\[
 (X_0,Y_0,P,p)=(V_{\rm reg}^{-1},Z,P,R/P),\qquad
 V_{\rm pt}=\eta V_{\rm reg},\qquad Q=P^3p.
\]
Here
\[
 f_0=(p-1)^{12}c_{12},\qquad
 c_{12}=\lambda_0c_q\eta^6/(u_s^{12}u_t^3),\qquad
 \eta=a_T\Theta'_Q(1)\Theta_Q(P^2)/\Theta_Q(P)^2.
\]
All these analytic units and their inverses are retained on the
zero-end domain. The original source functions are
\(s=P(p-1)/u_s\), \(t=Pu_t\), and \(q_{\rm old}=pPc_q\).

The zero period basis is \(L_0=B^{-1}\), \(T_0=hB^{-1}\):
\[
 L_0:(X_0,Y_0)\mapsto(P^{10}p f_0X_0,P^3Y_0),\qquad
 T_0:(X_0,Y_0)\mapsto(P^8p f_0X_0,P^6pY_0).
 \tag{I2ZC.1}
\]
Both keep \(P,p\). Their inverse functions use the reciprocal
coefficients at this same base point. In the common torus
\((X_c,Z_c)=(V_{\rm reg}/Z,Z)\), the complete periods are
\[
 L_c=(P^{23}p^{-1}f_0^2X_c,P^3p^3Z_c),\qquad
 T_c=(P^{11}p^{-1}f_0X_c,P^3p^2Z_c).
\]
For a saved common-parent word \(L_c^iT_c^j\), its full map is
\[
 \begin{aligned}
 X_c&=P^{23i+11j}p^{-i-j}f_0^{2i+j}X_0^{-1}Y_0^{-1},\\
 Z_c&=P^{3i+3j}p^{3i+2j}Y_0.
 \end{aligned} \tag{I2ZC.2}
\]
Thus the integral four-character matrix is
\[
 F_{i,j}=\begin{pmatrix}
 -1&-1&23i+11j&-i-j\\
 0&1&3i+3j&3i+2j\\
 0&0&1&0\\0&0&0&1
 \end{pmatrix},\qquad\det F_{i,j}=-1,
 \tag{I2ZC.3}
\]
with full \(f_0\)-column \((2i+j,0,0,0)^t\). The inverse is
the literal inverse matrix and reciprocal transported coefficient;
both are exported. On a generic ray \((x,y,h)\), the native ray is
\[
 (-x-y+(26i+14j)h,\ y-(3i+3j)h,\ h).
 \tag{I2ZC.4}
\]
The original physical vector \((-1,0,0,0)\) maps to
\((1,0,0,0)\). The genuine point coefficient acquires exactly
\(P^{26i+14j}p^{2i+j}f_0^{2i+j}\); its inverse retains every
power. In particular the p-character in this formula is carried
through the zero-end collar as a divisor character.

\par

\medskip\noindent\textbf{I2ZC.2. The relative four-cone construction}

\par\nobreak

For every saved common cone, I2ZC.4 and its recorded parent word
give a primitive regular native three-cone. The inverse of its
native upper-parent matrix sends all three columns to nonnegative
integers. The certificate checks uniqueness and the full rational
area of every one of the 72 native upper triangles.

The existing 420-chart zero parent already includes both actual
source bands and all periodic source-cut subdivisions. Its rays
are in \((\operatorname{ord}X_0,\operatorname{ord}Y_0,
\operatorname{ord}P,\operatorname{ord}p)\). The new relative
fan is constructed inside each such parent as follows. An upper
three-face is replaced by its pulled-back common triangles and
joined to the same actual lower ray. An upper two-face is replaced
by every common edge it contains and joined to its same two actual
lower rays. Parents with upper dimension at most one retain their
original rays. All lower rays are those of the executed source-cut
parent, including their nonzero original P-heights.

For each resulting four-cone the certificate computes its full
matrix \(N\), its integral inverse, and its determinant. It orders
the columns with determinant one. With parent matrix \(N_0\),
the coefficient map is literally
\[
 z_{0,a}=\prod_b z_b^{(N_0^{-1}N)_{ab}}.
 \tag{I2ZC.5}
\]
Every exponent is a nonnegative integer. On a common-face
localization the inverse is the matrix \(N^{-1}N_0\); all
coordinates that occur with negative exponents there are
invertible on that localization. The records retain both matrices.
These maps prove the asserted refinement to this actual regular
420-chart parent.

For each parent, the sums of the child slice volumes agree exactly.
For a regular cone on rays \(v_a\), its volume on
\(\operatorname{ord}P+\operatorname{ord}p=1\) is
\[
 \frac1{6\prod_a(v_{a,P}+v_{a,p})}.
 \tag{I2ZC.6}
\]
The common-face incidence calculation checks that every internal
three-face occurs twice and every boundary three-face once. The
generic boundary is exactly the entire pulled-back common fan.
The two original lower facets remain. These support and incidence
equalities give the proper chart gluing, with no additional face
selected by its characteristic value.

\par

\medskip\noindent\textbf{I2ZC.3. Full periodic source and line maps}

\par\nobreak

The four-dimensional valuation maps of I2ZC.1 are
\[
 H_L=\begin{pmatrix}1&0&10&1\\0&1&3&0\\0&0&1&0\\0&0&0&1\end{pmatrix},
 \quad
 H_T=\begin{pmatrix}1&0&8&1\\0&1&6&1\\0&0&1&0\\0&0&0&1\end{pmatrix}.
\]
The paired parent boundaries are those for \(L_0,T_0,L_0T_0^{-1}\),
with their actual full source and target equations. For a paired
source cone \(N_s\), target cone \(N_t\), and saved word
\((a,b)\), the complete map is
\[
 z'_\alpha=
 f_0^{[N_t^{-1}(a+b,0,0,0)^t]_\alpha}
 \prod_\beta z_\beta^{[N_t^{-1}H_L^aH_T^bN_s]_{\alpha\beta}}.
 \tag{I2ZC.7}
\]
Its inverse uses \(N_s^{-1}H_T^{-b}H_L^{-a}N_t\) and the
coefficient \(N_s^{-1}(-a-b,0,0,0)^t\). Since the periods fix
\(P,p\), these inverse unit columns cancel pointwise. The
certificate checks both localized exponent maps and the complete
unit inverse for every paired face, and accounts for the entire
remaining boundary by these periods and the two lower facets.

Every child retains its parent\textquotesingle s actual source band \(j=1\) or
\(j=2\). The original source functions on that band are
\[
 x_j=Q^jY_0,\qquad y_j=Q^{1-j}Y_0^{-1},\qquad x_jy_j=P^3p.
 \tag{I2ZC.8}
\]
Their exponent rows on a native ray \((x,y,r,z)\) are
\(y+3jr+jz\) and \(-y+3(1-j)r+(1-j)z\). Both rows are
nonnegative on every child. Their sum is exactly \(3r+z\).
The original analytic functions \(s,t,q_{\rm old}\) above are
then carried in these same charts and period squares.

The marked point graph remains the literal graph of
\((x_j,P^2)\), with its original two point-module generators
\((a-1)^{-1}\), \((c-1)^{-1}\), and the zero generator
\(x_j-1\). Its coefficient formulas are the actual I2P maps:
\[
 V_{\rm reg}=-\eta^{-1}P^{-2j}EA
 =\eta^{-1}P^{-2j}c^{-1}EC,\quad C=-cA,
\]
\[
 V_{\rm reg}=\eta^{-1}P^{-2j-2}(x_j-P^2)E_OC_O.
 \tag{I2ZC.9}
\]
Their pullbacks by I2ZC.5 are retained in every child record.
The finite marked comparison graphs use these actual generators
and the corresponding inverses; the full coefficient \(\eta\)
includes its theta derivative. The graph and coefficient arrows
give the original point-line boundary map used in the later
compact comparison.

\par

\medskip\noindent\textbf{I2ZC.4. Original determinant and exact scope}

\par\nobreak

The original four-form in the zero marking is
\[
 \Omega_{\rm orig}=
 \frac{2J_{ts}}{a_T}P\,dP\wedge dp\wedge
 d\log X_0\wedge d\log Y_0.
 \tag{I2ZC.10}
\]
The displayed sign and full nonzero unit are retained. On a ray
\((x,y,r,z)\), its canonical order is \(2r+z-1\). This is
recorded on every child together with the physical weights
\(N^{-1}(-1,0,0,0)^t\). The full derivatives of the base units
in I2ZC.2 and I2ZC.7 wedge to zero against the two base
differentials; their nonzero multiplicative units remain in the
actual chart and normal maps. No Spin extension is deduced by
discarding a canonical divisor.

The companion exports the complete native generic pullbacks,
four-cones, parent/inverse matrices, source rows, point/O modules,
period words and full coefficient columns. Input hashes are
verified before and after execution. Its exact counts and pass
status are recorded in its output; the earlier 72-triangle
preflight remains separately labelled as a preflight. The fixed
normal and moving-cap calculations use the actual new output and
are subsequent constructions.


% End expanded verified cp36 input


% Expanded verified cp36 input: notes/part3_e6_selected_i3_m2_return_support_09.tex


\par\bigskip\noindent\textbf{Preserved alternate construction: complete original auxiliary support 9}\par
\par\bigskip\noindent\textbf{P3K.111 supporting proof 9: Complete infinity and one lifts}\par\nobreak
\par

\medskip\noindent\textbf{Selected-I3 \(M_{2}\) N12: full infinity and one-end lifts}

\par\nobreak

\emph{Historical construction record.} All pending-execution statements in this supporting proof describe its preparation stage. The completed 145160-chart infinity and 149774-chart one-end execution, including the full-object comparison, is recorded in \path{records/e6_selected_i3_m2_corner_complete_equivalence.json}. The completed common input, compact coverage and return are I2F, CFL and ASI2 below.

Execution status: the exact common input is still running. The complete
coordinate construction below is prepared for that input; its final counts,
incidence gate and executed certificate are pending. No characteristic
pair or completed global return is assigned.

We retain every ray and full function of I2F.1-\/-10.
The generic boundary is the entire I2F common regular fan.
This construction extends that exact boundary across
the infinity and one endpoints. The zero endpoint
is the independently constructed I2ZC.1-\/-4 atlas.
All three use the original \(M_{2}\) point line and
translation by minus three G. Point-section graphs
and the compact return ledger remain further work.

\par

\medskip\noindent\textbf{I2L.1. The literal generic boundary at each endpoint}

\par\nobreak

At infinity the native characters are already
(X,Z,R) with \(X=X_c,Z=Z_c\). Attach the ray
\(A_{0}=(0,0,0,1)\) to every common three-cone, in the
ordered lattice (ord X,ord Z,ord R,ord u).
The full parameter arrow is
\[
 (X_c,Z_c,P,p)=(X,Z,Ru,u^{-1}),\qquad
 f_0=u^{-12}f_\infty.
 \tag{I2L.1}
\]
Its inverse is \(R=pP,u=p^{-1}\), with X,Z fixed.
All these equalities are the entire holomorphic
functions on their stated parameter overlaps.

At one, take a common cone with recorded one-parent
word(i,j). Its common ray(a,b,h) has native ray
\[
 \bigl(-b+(3i+3j)h,\ a-(23i+11j)h,\ h\bigr)
 \tag{I2L.2}
\]
in (ord \(X_{1}\),ord \(Y_{1}\),ord P). The common-input gate checks that these cones are
distinct, have total area thirty-six, and have
precisely the edge incidence of the native
parallelogram I2O.1 at z zero. The common period
edge identifications I2F.4 give the same
subdivisions from the two incident parent
triangles. Thus I2L.2 unwraps the actual common
quotient into the complete native fundamental
parallelogram. Join each cone to \(E_{0}=(0,0,0,1)\).

The full one-end character arrow, including the
additional \(v=p-1\) powers, is
\[
 \begin{pmatrix}\log X_c\\\log Z_c\\\log P\\\log v\end{pmatrix}
 =
 \begin{pmatrix}
 0&1&23i+11j&12(2i+j)\\
 -1&0&3i+3j&0\\
 0&0&1&0\\0&0&0&1
 \end{pmatrix}
 \begin{pmatrix}\log X_1\\\log Y_1\\\log P\\\log v\end{pmatrix}
 +u_p\log p+u_1\log c_{12},
 \tag{I2L.3}
\]
where
\(u_p=(-i-j,3i+2j,0,0)^t\) and
\(u_1=(2i+j,0,0,0)^t\).
This logarithmic notation denotes the literal
integer monomial and unit product; it introduces
no logarithm choice. The inverse character
matrix is integral. Every shared edge retains
all its incident common-owner words. Differences
of those words give exactly the full I2F.1
period arrows on the common open part.

\par

\medskip\noindent\textbf{I2L.2. The finite lower joins and their subdivisions}

\par\nobreak

At infinity, extract the common boundary intervals
on \(y+3h=0\) and on \(3x-11y=0\). Join each interval to
\(A_{0},A_{1}=(0,-1,0,1)\). On the first edge the joins
have determinant one; on the second they have
determinant eleven. Their support is precisely the
lower part of I2I.1, adjoining the upper stars
of I2L.1. Every join lies in one of the original
thirteen lower parents of I2I.2, and its parent is
determined by testing its four rays against
that parent\textquotesingle s inverse ray matrix.

At one, extract every upper boundary interval on
\(3y-x-36h=0\). For each \(j=0,...,11\) join it to
\(E_j=(0,j,0,1)\) and \(E(j+1)\). Every such join has
determinant three and lies in the actual parent
I2O.2 indexed by j. They cover the extrusion of
the native top edge through the full twelve-unit
v height. Together with the upper stars they
have exactly the support I2O.1.

For each singular lower cone, enumerate the
integral points of its half-open fundamental
parallelepiped, divide each nonzero vector by
its gcd, and insert the chosen primitive ray
in every containing lower cone. Each positive
barycentric coefficient is strictly less than
one on every changed cone; every child has
positive determinant strictly smaller than its
parent. All coefficients, old cones and children
are retained in chronological order.

Every new infinity ray satisfies either
\(y<-3h\) or \(3x-11y>0\). Every new one-end ray satisfies
\(3y-x>36h\). These strict inequalities put them
outside every upper star, whose corresponding
inequalities are weak in the opposite direction.
The generic boundary is therefore unchanged by
all these lower subdivisions.

The executed counts of lower stellar insertions and regular four-cones
will be inserted after both complete endpoint certificates pass. All ordered ray
matrices have determinant one and all cones
are distinct. On each original raw parent,
the sum of its child simplex volumes in
\(h+t=1\) is exactly its volume, using
\[
 \operatorname{vol}(B)
 =\frac{\det B}{6\prod_{v\text{ a ray of }B}(v_h+v_t)}.
 \tag{I2L.4}
\]
The explicit stellar operations establish
non-overlap and support; I2L.4 independently
checks the complete coverage of every parent.

\par

\medskip\noindent\textbf{I2L.3. Full affine morphisms to the original raw parents}

\par\nobreak

Write B for a final ordered four-ray matrix.
The child chart is the actual monomial chart
\((X,Z,R,u)=\theta^B\) at infinity, or
\((X_{1},Y_{1},P,v)=\theta^B\) at one. Its inverse ray
matrix has integral entries. For the raw
parent retain every primitive dual ray and
every integral half-open-parallelepiped
representative listed in I2I.2 or I2O.2.
They generate its entire dual semigroup by
subtracting the four coefficient floors from
any lattice point in the dual cone.

For each retained dual-semigroup generator m,
its exact pullback is \(\theta^{mB}\). The
certificate verifies every exponent is an
integer at least zero. Equal character sums
have equal products after this substitution,
so every binomial relation of the parent
semigroup is preserved. This is the full
affine morphism to the original parent,
including each determinant-eleven or
determinant-three singular parent. The
identity of characters gives the overlap
localizations and their inverse maps.

These finite subdivisions cover each raw
parent support; the resulting toric morphisms
are proper. This identifies their precise
targets. The earlier 149-cone and 168-cone
resolutions retain their own recorded maps
to the same raw parents. The construction
here does not use an unverified affine
morphism between the two regular refinements.

\par

\medskip\noindent\textbf{I2L.4. Every period face and its complete inverse}

\par\nobreak

At infinity use the actual lattice matrices
L and T of I2I.3, with full coefficient
\(f_{\rm infinity}^{2i+j}\) in the X coordinate for
the word \(L^iT^j\). The entire opposite-face
pairings are
\[
 L:\ x+23h=0\longrightarrow x=0,\qquad
 T:\ y+3h+t=0\longrightarrow y=0,
\]
\[
 L/T:\ 3x-11y+36h=0\longrightarrow3x-11y-11t=0.
 \tag{I2L.5}
\]
At one use I2O.3. The two pairings are
\[
 L:\ x=0\longrightarrow x-3h=0,\qquad
 T:\ 3y-x=0\longrightarrow3y-x-36h-36t=0.
 \tag{I2L.6}
\]
Their complete character-unit columns for
\(L^iT^j\) are \(p^{i-j,i,0,0}\) and
\(c_{12}^{0,j,0,0}\).

For each paired face, the record specifies
the unique adjacent source and target charts
Bs,Bt, the actual lattice map N, and the
integer matrix \(E=B_{t}^{-1} N B_{s}\). Every column
belonging to a shared face ray is exactly
the corresponding target coordinate vector.
Thus all negative entries occur only in
coordinates inverted on the stated overlap.
The inverse \(E^{-1}\) is integral and has the
same property in the reverse direction.
For every full character-unit column u the
actual unit exponent vector is \(B_{t}^{-1} u\);
the inverse is \(-B_{s}^{-1} N^{-1} u\). Multiplication
by \(E^{-1}\) verifies their complete cancellation,
including all parameter-dependent units.

The complete internal, generic and period-facet counts are pending
the executed arrays. Every internal
facet belongs to two cones. Every other
facet is either one of the prescribed
generic cones or one of I2L.5-\/-6; there is
no unpaired artificial boundary. The final incidence and opposite-face checks certify these
literal localized morphisms on the two complete atlases.

\par

\medskip\noindent\textbf{I2L.5. Source, point-line, physical weights and orientation}

\par\nobreak

At infinity the original source coordinates
are \(Z_{0}=R^{3}u^{2}Z,W_{0}=Z^{-1}\). Their exponent rows
on a ray (x,y,h,t) are y+3h+2t and minus y.
At one they are \(Z_{0}=pP^{3}/X_{1},W_{0}=X_{1}\), with
rows \(3h-x\) and x and the retained p unit
in Z0. Every displayed row is nonnegative
on every complete cone. Their products
are the actual \(Q=R^{3}u^{2}\) or \(Q=pP^{3}\),
respectively. The full inverse Tate map
of E81 gives the original Weierstrass
source in these charts.

The point-line maps are exactly
\[
 v_2=\frac{\eta P^2XZ}{f_{P^2}(R^3u^2Z)}
 \quad\hbox{and}\quad
 v_2=\frac{\eta P^2Y_1}{X_1 f_{P^2}(pP^3/X_1)}.
 \tag{I2L.7}
\]
They retain \(\eta=a_T*\Theta'_Q(1)*\Theta_Q(P^{2})/\Theta_Q(P)^{2}\),
its reciprocal and central value minus one, both point/O
Cartier generators and the reciprocal
arrow. Their normalized section graphs
will supply the subsequent fixed-normal
construction.

The physical cocharacters are (1,0,0,0)
at infinity and (0,1,0,0) at one. All
chart weights are the exact integer
vectors \(B^{-1} \delta\). The original volume is
\[
 \Omega=\frac{2J_{ts}}{a_T}\,
 dP\wedge dR\wedge d\log Z\wedge d\log V_{2T}.
 \tag{I2L.8}
\]
At infinity X=\(V_{2,T}\) and \(P=Ru\). The base
identity \(\mathrm dP\wedge \mathrm dR=-R \mathrm dR\wedge \mathrm du\)
and the character swap each contribute
a minus sign, giving
\[
 \Omega=+\frac{2J_{ts}}{a_T}R\,
 dR\wedge du\wedge d\log X\wedge d\log Z.
 \tag{I2L.9}
\]
At one \(R=(1+v)P\), \(Z=X_{1}^{-1},V_{2,T}=Y_{1}\);
the character map has determinant minus
one and the base Jacobian is P. Thus
\[
 \Omega=-\frac{2J_{ts}}{a_T}P\,
 dP\wedge dv\wedge d\log X_1\wedge d\log Y_1.
 \tag{I2L.10}
\]
The coefficient \(2J_{\rm ts}/a_T\) is the entire
nonvanishing source unit. Both canonical
orders are exactly \(2h+t-1\), retained on
every ray. The signs are computed from the displayed base and character
Jacobians of this original selected-I3 source.

The prepared executable is
\path{code/part3_root/e6_selected_i3_m2_corner_lifts.py}.
It records both full atlases, every raw-parent semigroup map,
every period face map and inverse, all full coefficient columns,
physical weights and original canonical orders. The numerical
certificate, source locks and counts will be recorded after the
common fan closes and the two complete lifts execute. The point
and O graph attachment and compact normal ledger remain subsequent
steps of the selected-I3 \(M_{2}\) calculation.


% End expanded verified cp36 input


% Expanded verified cp36 input: notes/part3_e6_selected_i3_m2_return_support_10.tex


\par\bigskip\noindent\textbf{Preserved alternate construction: complete original auxiliary support 10}\par
\par\bigskip\noindent\textbf{P3K.111 supporting proof 10: Generic curve normal maps}\par\nobreak
\par

\medskip\noindent\textbf{Selected-I3 \(M_{2}\): actual generic fixed curves and all three normal clutches}

\par\nobreak

This proof concerns the generic one-dimensional fixed components in the
executed N12 three-end completion. It uses the full 136,442-cone common
fan, the actual zero common lift and the root\textquotesingle s one/infinity lifts. The
companion records every curve separately, with its three original
character morphisms, their inverses, all base units, the point-line
arrow and the original canonical character. It does not assign a
global return or include any moving, lower-reference or compact
surface contribution.

\par

\medskip\noindent\textbf{I2CV.1. Components and actual end charts}

\par\nobreak

Let N be the matrix whose columns are an ordered regular common
three-cone. The physical vector in the original common lattice
\((X_c,Z_c,P)\) is \((1,0,0)\). The normal weights are therefore exactly
\(w=N^{-1}(1,0,0)\). A cone with all three weights nonzero gives a fixed
curve. The executed minimal-face/period inventory supplies its
component identifier, a single common chart occurrence and trivial
period stabilizer. These assertions are checked anew for every
component by the companion, against the complete fixed-face output.

The base is the actual exceptional parameter sphere with coordinate p,
and its three specified discs have coordinates p, \(v=p-1\) and \(u=p^{-1}\).
For every common cone the three completed endpoint atlases each give
one four-cone owner. Three of its columns lie over height zero for the
disc coordinate. Its remaining column is the actual transverse ray \(l_e\).
The companion checks that its base multiplicity is one. In particular,
the zero source-band subdivision can give \(l_e\) positive original
P-height. This column is retained in the calculation below.

Write \(B_e\) for the full endpoint ray matrix, \(F_e\) for its exact three-row
character map to the common chart, and

\[
                         E_e=N^{-1}F_eB_e.
\tag{I2CV.1}
\]

The three generic columns of \(E_e\) are the actual permutation matrix
identifying the common normals. The fourth column is denoted \(c_e\).
Thus, in the full chart arrows, the common normal coefficients are
the endpoint coefficients multiplied by the base monomials with
orders \(c_e\) and the units listed in I2CV.2. The physical derivative
square \(E_e w_e=w\) is checked including the zero weight of the base
column. No normal line is identified merely by its weight.

\par

\medskip\noindent\textbf{I2CV.2. Every period word, unit and inverse}

\par\nobreak

For the zero end let (i,j) be its recorded common L,T translation.
In native coordinates \((X_0,Y_0,P,p)\) its complete character matrix is

\[
 F_0=\begin{pmatrix}
 -1&-1&23i+11j&-i-j\\
 0&1&3i+3j&3i+2j\\
 0&0&1&0
 \end{pmatrix},
\tag{I2CV.2}
\]

and the full additional coefficient is \(f_0^{2i+j}\) in its first
component. Here
\[
 f_0=(p-1)^{12}\lambda_0c_q\eta^6/(u_s^{12}u_t^3),\qquad
 \eta=a_T\Theta'_Q(1)\Theta_Q(P^2)/\Theta_Q(P)^2.
\tag{I2CV.3}
\]
All of these functions are those of the actual SI2 source. The
inverse full four-character matrix and its inverse \(f_0\) powers are
stored by the zero lift and carried verbatim into each curve record.
Both actual source bands of that lift remain recorded.

For the one end, in \((X_1,Y_1,P,v)\), the full character matrix is
\[
 F_1=\begin{pmatrix}
 0&1&23i+11j&12(2i+j)\\
 -1&0&3i+3j&0\\
 0&0&1&0
 \end{pmatrix}.
\tag{I2CV.4}
\]
The remaining units have exponent vectors
\((-i-j,3i+2j,0)\) for p and \((2i+j,0,0)\) for \(f_1\),
with \(f_0=v^{12}f_1\). The full four-character inverse has rows
\[
 \begin{pmatrix}
 0&-1&3i+3j&0\\
 1&0&-23i-11j&-12(2i+j)\\
 0&0&1&0\\0&0&0&1
 \end{pmatrix},
\tag{I2CV.5}
\]
and inverse unit vectors \((3i+2j,i+j,0,0)\) for p and
\((0,-2i-j,0,0)\) for \(f_1\). Multiplication checks both the
character identity and the cancellation of each entire unit vector.

At infinity the native height is R and \(P=Ru\). The full map is
\[
 (X_c,Z_c,P)=(X_\infty,Z_\infty,Ru),\qquad
 F_\infty=\begin{pmatrix}1&0&0&0\\0&1&0&0\\0&0&1&1\end{pmatrix}.
\tag{I2CV.6}
\]
Its inverse is \(R=pP\), \(u=p^{-1}\). Thus the extra u in the original
P-height is retained in every normal order. All three maps are
actual holomorphic monomials in their recorded regular chart,
with the displayed extending units; the complete parent semigroup
checks were made in the endpoint construction.

The point-line coefficient is also carried through every map. In
the actual source band b,
\[
 v_2=\eta P^{2b}V_{\rm reg}/f_{P^2}(Q^bZ).
\tag{I2CV.7}
\]
The I2P point and origin module frames extend this equation at its
zeros and poles. In particular the extra origin factor \(P^{-2}\)
is retained. The source-period coefficient of (i,j) is
\(P^{3i+3j}p^{3i+2j}\); the point coefficient in comparison to common
band one is
\[
 P^{26i+14j+2-2b}p^{2i+j}f_0^{2i+j}
 \frac{f_{P^2}(Q^bZ_{\rm native})}{f_{P^2}(QZ_{\rm common})}.
\tag{I2CV.8}
\]
Its inverse is the reciprocal at the literal inverse source
coordinate. Both point/O module generators and the selected-height
unit remain in the chart-to-boundary arrow. The companion checks
the whole period vectors before using their normal derivatives.

\par

\medskip\noindent\textbf{I2CV.3. Signed line degrees and original determinant}

\par\nobreak

The direction of I2CV.1 is endpoint coefficient to common
coefficient. Consequently the three actual normal line degrees on
the parameter sphere are
\[
                              d_i=-\sum_{e=0,1,\infty}c_{e,i}.
\tag{I2CV.9}
\]
This is the divisor calculation in the three explicit parameter
discs: a common coefficient equal to \(t_e^{c_{e,i}}\) times the
regular endpoint coefficient has frame transition with the
reciprocal power. The full unit factors contribute through the
same boundary maps. They extend on their respective discs with
their fixed nonzero constants, so their exact based logarithms
give degree zero there. They are retained in each exported map,
including the p and \(f_1\) factors at the one end.

Let \(h_i\) be the common primitive heights and \(r_e\) the original
native height of the actual lower ray. The full character arrows
give the identity
\[
                     \sum_i h_i d_i=-1-\sum_e r_e.
\tag{I2CV.10}
\]
The original four-form has its original base-blowup factor. In
these charts its generic order is \(2h_i-1\), and the lower-ray
contribution is \(2r_e\). Therefore its restriction to the completed
curve is exactly
\[
 \sum_i(2h_i-1)d_i+2\sum_e r_e=-2-\sum_i d_i.
\tag{I2CV.11}
\]
The right side is the actual adjunction degree on this parameter
sphere. The companion checks this equality independently from
the full normal monomial maps. This preserves the original
determinant map and every positive lower source height.

\par

\medskip\noindent\textbf{I2CV.4. Scoped characteristic residues}

\par\nobreak

The fixed curve is the actual parameter sphere and its three
complex normal lines have the weights \(w_i\) and degrees \(d_i\) above.
Apply the orbifold Stokes expression SZ.57--59 / IE.16 with the
physical line \(M=U^{-1}\). Define
\[
 W=\prod_i w_i,\quad S=\sum_i d_i/w_i,\quad
 Q_w=\sum_iw_i^2,\quad T=\sum_iw_id_i.
\]
Expanding the stated universal expression to first order on the
curve gives its physical M contribution
\[
               k_M=S/W,\qquad \ell_M=(Q_wS-2T)/W,
 \qquad A_M=(4k_M-\ell_M)/24.
\tag{I2CV.12}
\]
The companion independently expands the universal-U numerator
and verifies the Hermitian-dual sign in I2CV.12 for every curve.
All rational arithmetic is exact. It records each triple of
weights and degrees, all three endpoint maps and inverses, the
full point/source coefficient vectors and the original volume
identity before summing this scoped curve contribution.

The final machine result is source and input locked at
\path{code/part3/e6_selected_i3_m2_fixed_curve_normals.json}. It includes
all generic curve components certified by the complete face
inventory. Compact surfaces, endpoint interiors, lower reference
doubles, moving caps and the outer physical transfer remain
separately recorded constructions.


% End expanded verified cp36 input


% Expanded verified cp36 input: notes/part3_e6_selected_i3_m2_return_support_11.tex


\par\bigskip\noindent\textbf{Preserved alternate construction: complete original auxiliary support 11}\par
\par\bigskip\noindent\textbf{P3K.111 supporting proof 11: Whole compact surface/curve/corner ledger}\par\nobreak
\par

\medskip\noindent\textbf{Selected-I3 \(M_{2}\): complete compact fixed-component ledger}

\par\nobreak

This proof concerns the actual N12 periodic completion of SI2, with
\(M_2=\mathcal O([2G]-[O])\), translation by \(-3G\), the original
selected height \(q=t-t_c(s)\), and the cubic good cap and four-root-turn
reference retained. Its three endpoint atlases are I2ZC and I2L, glued
by the entire I2F common fan. The full characteristic contribution
of their compact generic and corner-interior fixed components is

\[
                  (k_M,\ell_M)_{\mathrm{compact}}=(2723,19124),
 \qquad A_{M,\mathrm{compact}}=-343.
 \tag{CFL.1}
\]

Here \(M=U^{-1}\) is the original physical nonzero-vector line.
The twelve lower/reference surfaces, the selected moving cap with
its whole-reference fan comparison, and the original outer
parameter transfer are the remaining pieces of the full boundary
ledger. CFL.1 is the contribution of the compact strata specified
below.

\par

\medskip\noindent\textbf{CFL.1. The actual three-ended family and fixed components}

\par\nobreak

Retain the complete source functions

\[
 Q=P^2R,\quad R-P=s u_s,\quad t=P u_t,\quad q=R c_q,
 \quad p=R/P,\quad
 \eta=a_T\Theta'_Q(1)\frac{\Theta_Q(P^2)}{\Theta_Q(P)^2},
\]
\[
 c_{12}=\frac{\lambda_0c_q\eta^6}{u_s^{12}u_t^3},
 \qquad f_0=(p-1)^{12}c_{12}.
 \tag{CFL.2}
\]

The full \(\eta\) and its reciprocal remain in the coefficients;
its central value is \(-1\), while \(a_T(0)=1\). On the common
torus set
\((X_c,Z_c,P)=(V_{\mathrm{reg}}/Z,Z,P)\), with
\(V_{\mathrm{pt}}=\eta V_{\mathrm{reg}}\). Its actual periods are

\[
 L_c:(X_c,Z_c,P)\longmapsto
 (P^{23}p^{-1}f_0^2X_c,P^3p^3Z_c,P),
\]
\[
 T_c:(X_c,Z_c,P)\longmapsto
 (P^{11}p^{-1}f_0X_c,P^3p^2Z_c,P).
 \tag{CFL.3}
\]

Both fix all the functions in CFL.2. Thus the word \(L_c^iT_c^j\)
has full unit vectors
\((-i-j,3i+2j,0)\) for \(p\) and \((2i+j,0,0)\) for \(f_0\).
Its original lattice shift is
\(((23i+11j)h,(3i+3j)h,0)\).

I2F constructs the complete regular common fan from the three
specified parent decompositions. Its 136,442 cones cover the
original area -36 period domain; all 204,429 interior edges and
468 boundary edges have their recorded incidence and full period
arrows. The two opposite-face words generate the original period
lattice. Every common affine coordinate is the primitive monomial
of its ordered determinant-one ray matrix. The parent maps
I2F.6--7 include the entire \(p,f_0\) coefficients and their inverse
characters.

The zero, one and infinity families are the proper subdivisions
I2ZC and I2L of their stated raw endpoint parents. Their recorded
semigroup maps pull every parent generator back to a monomial
with nonnegative integer exponents. Equal sums of parent
characters have identical products after pullback, so all parent
relations hold. The support and facet checks give the exact
gluing to the common fan, with every other boundary accounted
for by an original period face or a labelled lower face.
The full zero construction retains both source bands. Infinity
retains \(P=Ru\), and one retains \(f_0=v^{12}c_{12}\), \(v=p-1\).

For an ordered local ray matrix \(B\), the actual circle weights
are \(B^{-1}\delta\). The vectors are
\(\delta_c=(1,0,0)\), \(\delta_0=(-1,0,0,0)\),
\(\delta_1=(0,1,0,0)\), and \(\delta_\infty=(1,0,0,0)\)
in their respective named lattices. A fixed coordinate stratum
sets precisely the nonzero-weight coordinates to zero. Its
normal face is their span. Localization along a common toric
face gives the exact overlap of these fixed strata. Across a
period facet the full chart map and its inverse identify the
face, its tangent coordinates and its normal derivative. The
complete fixed-face inventory performs these identifications,
preserving every period word and stabilizer.

The components used here have trivial period stabilizer, checked
from their full orbit cycles. There are 106,166 generic fixed
curves and 15,138 generic fixed surfaces. The corner-interior
components are listed in CFL.8. Their complete stars close in the
three executed atlases. The twelve one-end components with a
lower source boundary remain separately labelled in that inventory.

\par

\medskip\noindent\textbf{CFL.2. Every full endpoint normal morphism}

\par\nobreak

Let \(N\) be the ordered matrix of a common three-cone, and let
\(B_e\) be its actual four-cone owner at \(e=0,1,\infty\).
Its three generic columns have disc-coordinate height zero;
its fourth column is the original lower ray \(\lambda_e\),
of disc multiplicity one. The common-character map \(F_e\) gives

\[
                         E_e=N^{-1}F_e B_e.
 \tag{CFL.4}
\]

The three generic columns are the recorded coordinate permutation.
Write \(c_e\) for its lower column. The full local morphism is

\[
 z_i=\varepsilon_{e,i}t_e^{c_{e,i}}z_{e,i},\qquad
 z_{e,i}=\varepsilon_{e,i}^{-1}t_e^{-c_{e,i}}z_i,
 \tag{CFL.5}
\]

on its specified coefficient localization. The matrices and
units extend by the endpoint semigroup maps at the regular
coordinate boundary. The base-disc coordinates are \(p,v=p-1\),
and \(u=p^{-1}\).

For a zero owner with common word \((i,j)\),

\[
 F_0=
 \begin{pmatrix}
 -1&-1&23i+11j&-i-j\\
 0&1&3i+3j&3i+2j\\
 0&0&1&0
 \end{pmatrix},\qquad u_{f_0}=(2i+j,0,0)^t.
 \tag{CFL.6}
\]

Here \(p\) is a divisor coordinate; its powers remain in the
matrix. The full four-character inverse and reciprocal
\(f_0\)-column are recorded by I2ZC. Every owner retains the
actual lower ray of the source-band subdivision, including
all positive original \(P\) heights.

For a one owner,

\[
 F_1=
 \begin{pmatrix}
 0&1&23i+11j&12(2i+j)\\
 -1&0&3i+3j&0\\
 0&0&1&0
 \end{pmatrix}.
 \tag{CFL.7}
\]

Its unit vectors are \((-i-j,3i+2j,0)\) for \(p\) and
\((2i+j,0,0)\) for \(c_{12}\). Its full inverse character
matrix is

\[
 \begin{pmatrix}
 0&-1&3i+3j&0\\
 1&0&-23i-11j&-12(2i+j)\\
 0&0&1&0\\0&0&0&1
 \end{pmatrix},
 \tag{CFL.8}
\]

with inverse coefficient columns
\((3i+2j,i+j,0,0)\) and \((0,-2i-j,0,0)\).
Multiplication gives the identity and cancels both complete
unit columns. At infinity,

\[
 (X_c,Z_c,P)=(X_\infty,Z_\infty,Ru),\qquad
 F_\infty=
 \begin{pmatrix}
 1&0&0&0\\0&1&0&0\\0&0&1&1
 \end{pmatrix}.
 \tag{CFL.9}
\]

The inverse is \(R=pP,u=p^{-1}\). The \(u\) factor in the
original \(P\) is retained in the normal and determinant maps.

For a period facet with source and target matrices \(B_s,B_t\),
original lattice map \(H\), and full coefficient column \(u\),
the affine overlap is

\[
 z'_\alpha=
 \varepsilon^{(B_t^{-1}u)_\alpha}
 \prod_\beta z_\beta^{(B_t^{-1}HB_s)_{\alpha\beta}}.
 \tag{CFL.10}
\]

Its inverse uses \(B_s^{-1}H^{-1}B_t\) and
\(-B_s^{-1}H^{-1}u\). The recorded localizations make every
negative-power coordinate invertible. Composing both maps
gives the identity as the full analytic function. Thus CFL.5
is the normal derivative of the actual source/line chart map.

All coefficient functions in CFL.10 are invariant under the
physical circle. Along a fixed component their derivative on
a nonzero-weight normal direction vanishes by equivariance.
Differentiating the complete coordinate arrow therefore gives
its written coefficient restricted to that component, with
the recorded monomial normal factors. The zero-weight tangent
block and the signed normal eigenlines retain their ordering.

\par

\medskip\noindent\textbf{CFL.3. Original point modules and determinant}

\par\nobreak

The source band \(b\) has

\[
 x_b=Q^bZ,\qquad y_b=Q^{1-b}/Z,\qquad x_by_b=P^2R.
 \tag{CFL.11}
\]

The proper point graph is the blowup of the literal ideal
\((x_b,P^2)\), with charts

\[
 x_b=P^2a,\quad ay_b=R;\qquad
 P^2=x_bc,\quad y_b=Rc.
 \tag{CFL.12}
\]

Their overlap is \(c=a^{-1}\), and the point section is
\(a=1\) or \(c=1\). The second chart has the explicit
two-chart resolution of \(x_bc=P^2\) in I2P.2. Pulling
these equations back by the endpoint monomials gives the
finite marked source graphs and their original source maps.

Let \(E,E_O\) be the entire convergent theta-product units of
I2P.3 and I2P.7. With the point-divisor generators
\(e_a=(a-1)^{-1},e_c=(c-1)^{-1}\) and zero-divisor generator
\(e_O=x_b-1\), the full coefficient arrows are

\[
 V_{\rm reg}=-\eta^{-1}P^{-2b}EA
       =\eta^{-1}P^{-2b}c^{-1}EC,\qquad C=-cA,
\]
\[
 V_{\rm reg}=
       \eta^{-1}P^{-2b-2}(x_b-P^2)E_OC_O.
 \tag{CFL.13}
\]

Each inverse uses the reciprocal at the actual inverse source
point. The \(c^{-1}\) has its specified Cartier type, and
the extra \(P^{-2}\) at \(O\) is retained. The source period
has its simultaneous band change \(b\mapsto b-1\); substitution
leaves \(A,C,C_O\) unchanged. The full suspension is

\[
 (Z,V_{\rm reg})\longmapsto(P^{-3}Z,d_BV_{\rm reg}),
 \quad d_B=\lambda_0s^{12}q\eta^6t^{-3},
\]
\[
 v_2'=d_BP^{2(b'-b)}
       \frac{f_{P^2}(x_b)}{f_{P^2}(x'_{b'})}v_2,
 \quad x'_{b'}=Q^{b'-b}P^{-3}x_b.
 \tag{CFL.14}
\]

These are the literal SI2 Miller Hom and I2P source/Cartier
squares. For the endpoint-to-common word \((i,j)\), the
genuine point coefficient relative to common band one is

\[
 P^{26i+14j+2-2b}p^{2i+j}f_0^{2i+j}
 \frac{f_{P^2}(Q^bZ_{\rm native})}
      {f_{P^2}(QZ_{\rm common})}.
 \tag{CFL.15}
\]

Equations CFL.12--14 cancel its point and origin zeros and
poles in their specified modules. They give the actual
boundary line morphism with exactly the \(p,f_0,\eta\) and
selected-height functions of CFL.5--10. The normal calculation
here concerns the stated completed periodic atlas. Its
original point-line boundary identification is this finite
graph pullback, with its reciprocal.

The original volume forms are

\[
 \Omega_0=+\frac{2J_{ts}}{a_T}
       P\,dP\wedge dp\wedge d\log X_0\wedge d\log Y_0,
\]
\[
 \Omega_1=-\frac{2J_{ts}}{a_T}
       P\,dP\wedge dv\wedge d\log X_1\wedge d\log Y_1,
\]
\[
 \Omega_\infty=+\frac{2J_{ts}}{a_T}
       R\,dR\wedge du\wedge d\log X_\infty\wedge d\log Z_\infty.
 \tag{CFL.16}
\]

Here \(J_{ts}=\det\partial(t,s)/\partial(P,R)\); the signs
are the exact base and character Jacobians. A native ray
\((a,b,r,z)\) has original canonical order \(2r+z-1\);
a generic ray \((a,b,h)\) has order \(2h-1\).
The base-unit derivatives wedge to zero against the two
base differentials in CFL.16, while their nonzero
multiplicative units remain in every coordinate and normal map.

For every common coordinate line, the completed degree is

\[
                         d_i=-\sum_e c_{e,i}.
 \tag{CFL.17}
\]

This follows from the actual direction of CFL.5: the frame
transition is the reciprocal of the coefficient transition.
Each remaining unit has its actual logarithm on the endpoint
disc after fixing its nonzero constant. Its degree is zero,
and its full boundary value stays in the bundle arrow.
The independently checked original adjunction identity is

\[
 \sum_i(2h_i-1)d_i+2\sum_e r_e=-2-\sum_i d_i,
 \tag{CFL.18}
\]

where \(r_e\) is the original native height of the actual
lower ray. This includes every positive lower height and
the full infinity factor \(P=Ru\).

\par

\medskip\noindent\textbf{CFL.4. Generic fixed surfaces and saturated quotient fans}

\par\nobreak

Each of the 15,138 generic surfaces has precisely two common
chart occurrences, two ordered normal rays \(n_-,n_+\), and
weights \(w_-<0<w_+\). Every period cycle is trivial.
Align each occurrence by the inverse of its full recorded
period word. The remaining rays of the two aligned charts
are \(v_0,v_\infty\). Their quotient by the normal span is
the saturated rank-one lattice, because either complete
three-ray matrix is unimodular. The exact wall equation is

\[
 v_0+v_\infty+d_-^Fn_-+d_+^Fn_+=0.
 \tag{CFL.19}
\]

The two projected rays are opposite primitive generators.
The two actual affine fibre coordinates are reciprocal,
with the full base unit of CFL.10, and their union is the
\(\mathbb P^1\) fibre. The normal fibre degrees are
\(d_\pm^F\). Its two coordinate sections \(S_0,S_\infty\)
are the tangent-zero curves in these charts. Their
self-intersections in the surface and their two normal
degrees are the three entries of CFL.17, with the actual
recorded tangent and normal indices.

At an endpoint, choose a complete aligned chart
\(B=(n_-,n_+,t_0,t_1)\). The last two rows of \(B^{-1}\)
give the saturated quotient lattice \(\mathbb Z^2\).
Every projected tangent ray is integral and primitive,
and consecutive rays have determinant one. The whole
projected fan is a chain between two opposite boundary
section rays of base height zero. Every other ray has
positive base height. The interior valence is two and
the boundary valence one; the exact traversal visits
all charts and all rays.

For an internal tangent ray \(r\) with neighbours \(r_l,r_r\),
the complete four-dimensional relation is

\[
 r_l+r_r+d_-^r n_-+d_+^r n_+ + s_r r=0.
 \tag{CFL.20}
\]

The inverse of \((n_-,n_+,r,r_l)\) computes it; its fourth
independent coefficient is zero. In the two adjacent
charts it gives \(z'=z^{-1}\) and the normal coefficient
changes \(a_\pm'=z^{-d_\pm^r}a_\pm\), including the entire
units of CFL.10. Reciprocal frames prove the degrees
\(d_\pm^r\); the tangent transverse coefficient proves
the self-intersection \(s_r\). The full fibre is
\(\sum_r m_r C_r\), with \(m_r\) the actual base height.
Every relation \(m_rs_r+m_{r_l}+m_{r_r}=0\) is checked,
with boundary multiplicity zero. Thus the original
base-function divisor has zero intersection with each
vertical curve.

\par

\medskip\noindent\textbf{CFL.5. Proper blowdowns and integral intersection basis}

\par\nobreak

Each executed chain contains a multiplicity-one component.
The saved anchor is one of these. Each other ray is
contracted only after its saturated quotient vector is
the sum of its two neighbours. The actual proper
ordinary blowup has affine maps

\[
       (u,v)\mapsto(u,uv),\qquad
       (u',v')\mapsto(u'v',v'),
 \tag{CFL.21}
\]

with overlap \(u'=v^{-1},v'=uv\). The full parent
monomials and normal units compose with these maps.
The anchor remains, and the final three-ray half-plane
is the ordinary \(\mathbb P^1\)-fibre chart over the disc.
There are exactly 52 recorded contractions: 40 at zero,
none at one and 12 at infinity. Every removed ray and
its two then-current neighbours are saved.

After contraction, the fibre charts are reciprocal and
their base transition is \(z_\beta=g_{\alpha\beta}(p)z_\alpha\),
with the full nonzero function from CFL.10. Thus the
contracted surface is the actual
\(\mathbb P(\mathcal O\oplus L)\) over the parameter sphere,
with \(L\) defined by that cocycle and with the two
specified sections. CFL.21 reconstructs the executed
surface by the recorded blowups.

Use the integral divisor basis

\[
 \mathcal B=\bigl(S_0,F,\{C_{e,r}:r\ne r_{e,\rm anchor}\}\bigr).
 \tag{CFL.22}
\]

The multiplicity-one anchors express the omitted components
integrally by the full fibre relation. Conversely the
chronological maps CFL.21 express this basis in the section,
fibre and integral exceptional classes. Hence it is an
integral basis. Its matrix \(I\) is given by
\(S_0^2=s_0,\ S_0F=1,\ F^2=0\), the full self-intersections
CFL.20, the chain adjacency and the actual section incidences.
Different endpoints are disjoint. Its inverse is integral,
which is also checked directly.

With normal degrees \(d_\pm^0\) on \(S_0\), the actual
normal Chern vectors are

\[
 \mathbf n_\pm=I^{-1}
       (d_\pm^0,d_\pm^F,(d_\pm^{e,r})_{e,r})^t.
 \tag{CFL.23}
\]

The second section is reconstructed from its pairings
\((0,1,\text{its component incidences})\). Its computed
self-intersection and both normal pairings agree with
the independently obtained values from CFL.17.
This verifies the full three-end normal-line gluing.

Adjunction on the listed rational curves gives

\[
 \mathbf K=I^{-1}
       (-2-s_0,-2,(-2-s_{e,r})_{e,r})^t.
 \tag{CFL.24}
\]

For \(b\) recorded blowups, the explicit maps give
\(K^2=8-b,\ c_2=4+b,\ p_1(TF)=-3b\). The matrix
values in CFL.24 satisfy these identities for every
surface. The normal intersections are the actual
products \(\mathbf n_-^tI\mathbf n_-\),
\(\mathbf n_-^tI\mathbf n_+\), and
\(\mathbf n_+^tI\mathbf n_+\), with all section and
wall data retained.

\par

\medskip\noindent\textbf{CFL.6. Compact corner points, curves and surfaces}

\par\nobreak

Each corner-interior component is aligned by the inverse of
its entire recorded period word. The full affine maps
CFL.10 identify its tangent and normal coordinates.
All these components have trivial period stabilizer.

An isolated point has one complete regular four-chart
star. Its four physical weights are the exact entries
of \(B^{-1}\delta_e\). A corner curve has two complete
charts. Its normal rays \(n_1,n_2,n_3\) and tangent rays
\(v_0,v_\infty\) obey

\[
                  v_0+v_\infty+\sum_{i=1}^3d_i n_i=0.
 \tag{CFL.25}
\]

The tangent rays project to opposite primitive generators
of the saturated quotient. Their two affine coordinates
give the actual compact \(\mathbb P^1\), and the coefficient/
frame calculation of CFL.20 gives its normal degrees
\(d_i\), with all aligned period and unit factors.

There are 17 compact corner surfaces, all at zero.
Project their full aligned stars by the saturated
rank-two normal quotient. The projected rays in cyclic
order have determinant one for every consecutive pair,
including the last-to-first pair. The list of those
pairs is exactly the list of chart cones. Thus the
projected fan covers the plane and its ordinary
affine gluing gives the complete compact toric surface.

On these corner-interior components both original endpoint
base functions vanish. The full unit coefficients restrict
to their fixed nonzero constants. Aligning a chart by its
recorded period word carries those constants as part of the
coordinate map. The trivial period-cycle check says that
their product around every alignment cycle is the identity.
Thus the aligned coordinate and normal maps give the stated
toric surface and its actual equivariant normal bundles.
This proves the torus lift used in the following calculation.

For an original integral auxiliary cocharacter \(a\),
each chart has the exact weights \(B^{-1}a\).
Its two normal and two tangent entries give

\[
 \int_F c_1(N_i)c_1(N_j)=
 \sum_\sigma
 \frac{n_{\sigma,i}(a)n_{\sigma,j}(a)}
      {t_{\sigma,1}(a)t_{\sigma,2}(a)},\qquad
 \int_Fp_1(TF)=
 \sum_\sigma
 \frac{t_{\sigma,1}(a)^2+t_{\sigma,2}(a)^2}
      {t_{\sigma,1}(a)t_{\sigma,2}(a)}.
 \tag{CFL.26}
\]

The auxiliary vectors lie in the original four-lattice.
Two distinct choices among
\((1,11,103,1009),(2,17,89,701),(5,13,29,53),(7,31,113,809)\)
have nonzero tangent entries throughout each complete
star. Both computations give identical integral values.
The actual torus action on the aligned star gives these
normal lifts, with coordinate and period maps CFL.10.
The tangent number also satisfies \(p_1=12-3r\), where
\(r\) is the number of rays of the complete smooth fan.

The independent root replay constructs the whole cyclic
divisor intersection matrix of each of the 17 surfaces.
It retains every divisor and every principal-character
relation, represents each normal line by its actual
support-function divisor \(D_i=-\sum_r a_i(r)D_r\), and
evaluates the products in that full matrix. This
independently reproduces all three normal intersections
and the tangent \(p_1\) of CFL.26.

\par

\medskip\noindent\textbf{CFL.7. Physical orientation and weighted residues}

\par\nobreak

The chart morphisms are holomorphic with the original
complex Jacobians CFL.16. The physical circle is the
displayed primitive cocharacter throughout. Use the
actual circle-first quotient and line \(M=U^{-1}\)
of SZ.57--59 and IE.16. A finite-stabilizer quotient
chart has its actual integration factor \(1/d\).
These strata lie in the quotient interior. Removing
the invariant normal tube of each fixed component gives
the fixed-boundary residue of

\[
                         \frac{4x^3-p_1(TZ)x}{24}.
 \tag{CFL.27}
\]

The formulas below expand the original equivariant
Euler denominator for every recorded nonzero integer
normal weight.

A point with four weights \(w_i\) contributes

\[
 k_M=-\frac1{\prod_iw_i},\qquad
 \ell_M=-\frac{\sum_iw_i^2}{\prod_iw_i}.
 \tag{CFL.28}
\]

For a fixed curve set \(W=\prod_iw_i\),
\(S=\sum_id_i/w_i\), \(Q_w=\sum_iw_i^2\) and
\(T=\sum_iw_id_i\). Its pair is

\[
                  k_M=\frac SW,\qquad
                  \ell_M=\frac{Q_wS-2T}{W}.
 \tag{CFL.29}
\]

For a surface use weights \(a,b\), intersections
\(x=\int c_1(N_-)^2,\ y=\int c_1(N_-)c_1(N_+),\
z=\int c_1(N_+)^2\), and \(p=\int p_1(TF)\).
Set

\[
                   B_2=\frac{x}{a^2}+\frac{y}{ab}+\frac{z}{b^2}.
 \tag{CFL.30}
\]

Then

\[
 k_M=-\frac{B_2}{ab},\qquad
 \ell_M=-\frac{p-x-z-2(a/b+b/a)y+(a^2+b^2)B_2}{ab}.
 \tag{CFL.31}
\]

Indeed, the \(U\) Euler denominator is
\(\prod_i(c_1(N_i)+w_i u)\). Its reciprocal has the
factor \(1-S/u\) in complex dimension one, and

\[
 1-\frac{x_-/a+x_+/b}{u}
 +\frac{x_-^2/a^2+x_-x_+/(ab)+x_+^2/b^2}{u^2}
\]

in complex dimension two. The \(u^{-1}\) coefficients
of \(u^3\) and
\(u[p_1(TF)+\sum_i(x_i+w_i u)^2]\) give the \(U\)
pair. The literal physical dual changes both signs,
giving CFL.28--31. This uses the actual physical line
and all finite normal weights.

\par

\medskip\noindent\textbf{CFL.8. Complete compact sums and ownership}

\par\nobreak

The 106,166 generic curves have the full normal clutches
of I2CV and CFL.4--18. Their exact pair is

\[
 k_{\rm curves}=
 -\frac{5884517677661433720832516179200147039685023720717}
 {692517323612787928519526807587126340256000000},
\]
\[
 \ell_{\rm curves}=
 -\frac{83730187907160166050180995681772221650262376569}
 {53270563354829840655348215968240487712000000}.
 \tag{CFL.32}
\]

The 15,138 completed generic surfaces have pair

\[
 k_{\rm surfaces}=
 \frac{5943727908830327088720935721248846341776911720717}
 {692517323612787928519526807587126340256000000},
\]
\[
 \ell_{\rm surfaces}=
 \frac{551192698986630725261025427857074561444990376569}
 {53270563354829840655348215968240487712000000}.
 \tag{CFL.33}
\]

Thus

\[
                  (k_M,\ell_M)_{\rm generic}
                           =(171/2,35101/4).
 \tag{CFL.34}
\]

The complete corner-interior table is

\begin{longtable}[]{@{}lrrrrr@{}}
\toprule\noalign{}
End & Points & Curves & Surfaces & \(k_M\) & \(\ell_M\) \\
\midrule\noalign{}
\endhead
\bottomrule\noalign{}
\endlastfoot
zero & 6163 & 72 & 17 & \(889/18\) & \(-212705/36\) \\
one & 4764 & 2340 & 0 & \(2736\) & \(17092\) \\
infinity & 7304 & 701 & 0 & \(-1331/9\) & \(-7513/9\) \\
\end{longtable}

Their sum is

\[
                    (k_M,\ell_M)_{\rm corners}
                            =(5275/2,41395/4).
 \tag{CFL.35}
\]

Equations CFL.34--35 prove CFL.1. Both generic types
have already been completed across the three parameter
discs. The corner-interior tags refer to their disjoint
complete fixed stars, with ownership proved by the
full period-orbit inventory. Each component in the
table is counted once.

\par

\medskip\noindent\textbf{CFL.9. Exact verification and remaining boundary}

\par\nobreak

The complete curve certificate is
\path{code/part3/e6_selected_i3_m2_fixed_curve_normals.json},
SHA256
\path{219b4314608a7d0c94f8586624a6e9643eeb6324e40a197502bae6c095f6e62b},
with 7,294,790 checks. It records every full endpoint
and inverse map, actual lower ray, point-period
coefficient and original determinant identity.

The complete surface certificate is
\path{code/part3_root/e6_selected_i3_m2_surface_stars.json},
SHA256
\path{000e22cf37479c760c6d9a5708cfe0f36e6d8887d5ad37732649825e76f177c4},
with 2,378,991 checks. It retains every common section
clutch, aligned surface star, saturated quotient ray,
chronological blowdown, integral intersection matrix,
normal Chern vector and actual tangent number.

The complete corner certificate is
\path{code/part3_root/e6_selected_i3_m2_corner_normals.json},
SHA256
\path{c8de0d7e5ef18f5a65bf648629cc0c03260449fb19a597aa1fa61dd656437cbc},
with 50,430 checks. It retains every aligned star,
exact wall relation, complete quotient fan and both
auxiliary-cocharacter computations. All three
executions have empty error streams.

The independent root arithmetic replay is
\path{code/part3_root/e6_selected_i3_m2_compact_arithmetic_peer.json},
SHA256
\path{a3804a52bad8b979ec072b3e2722a1b1fee6ae22a398543fb58764cb36eb7cf6}.
Its 856,099 checks multiply the original equivariant
Euler factors independently for all components and
reconstruct all 17 corner surface divisor matrices.
Every component value and the exact totals agree.

Surface and corner loading uses the new read-only
filtered JSON reader. Every mathematical function AST
and the complete geometric-computation suffix of both
programs remain unchanged. Its independent certificate
passes 151 checks, including standard-JSON decoding
of complete actual records and exact retained global
identifiers. The predecessor files, large-load failures
and equivalence records remain preserved under
\path{records/attempts/e6_selected_i3_m2_normal_filtered_loaders/.}

The compact calculation retains the source and point-line
arrows CFL.11--15, original physical action and complete
determinant CFL.16--18. The full original return also
requires the twelve lower doubles, selected moving
component and same-fan reference, and the outer
source/root/annular/parameter boundary. The whole
scalar-reference correction RFB acts on that entire
reference word. CFL.1 supplies the compact fixed term
for the stated final attachment.

The source-locked companion
\path{code/part3/e6_selected_i3_m2_compact_ledger_checks.py}
independently verifies the full character/unit inverses,
original physical vectors, both point and origin module
directions, formal Euler-denominator expansions, exact
scoped sums and all cited large-result hashes. Its fresh
JSON passes 109 checks. The complete large-row replay
remains the separately identified root certificate above.


% End expanded verified cp36 input


% Expanded verified cp36 input: notes/part3_e6_selected_i3_m2_return_support_12.tex


\par\bigskip\noindent\textbf{Preserved alternate construction: I2P actual atlas input}\par
E81 and the locked common/end atlases I2F and its actual endpoint parents and lifts. It establishes

\par\bigskip\noindent\textbf{Preserved alternate construction: I2P exact exterior and lower-source graph scope}\par
On a common/end chart all x,y,P,R are the explicit monomials and
full units of I2F.8-\/-9, I2I.3, I2O.3 and the complete zero-parent source bands. Pulling I2P.1 back
means blowing up the actual ideal generated by those two functions,
not only a valuation vector. The source and line coefficient
morphisms above specify the marked finite graphs needed by that
pullback. Period-translated graph owners must be carried with the
full I2P.8-\/-9 maps when constructing the final compact normal ledger.
This last global assembly is not asserted by the present local note.



\par\bigskip\noindent\textbf{Preserved alternate construction: I2M retained infinity period locator}\par
arrows of I2I.1. For the translate

\par\bigskip\noindent\textbf{Preserved alternate construction: native-star input scope}\par
This proof uses the executed zero endpoint and its actual common lift,

\par\bigskip\noindent\textbf{Preserved alternate construction: old zero atlas and eleven-chart extraction}\par
They are the unchanged lower facets of the 420-chart source-compatible
parent. This follows both from the relative construction I2ZC.2 and
from the persisted 143,270-cone output. The companion extracts every
cone containing either physical face \(\langle A,B_1\rangle\) or
\(\langle B_2,B_3\rangle\); there are seven and four, respectively.
It checks every parent map and exports the original source band,
normal characters, original canonical orders and full period units.



\par\bigskip\noindent\textbf{Preserved alternate construction: native normal derivative chart count}\par
For every one of the eleven resolved moving-star charts, the
companion exports the actual parent matrix and its inverse. Align

\par\bigskip\noindent\textbf{Preserved alternate construction: old eleven-chart source fan and nine blowdowns}\par
The quotient tangent rays of the eleven charts, in the aligned
\((\operatorname{ord}Z,\operatorname{ord}P)\) lattice, are exactly
\[
 (1,0),(3,1),(2,1),(1,1),(0,1),(-1,1),(-2,1),(-3,1),
                 (-4,1),(-5,1),(-6,1),(-1,0).
 \tag{NS.23b}
\]
Every adjacent pair is a primitive cone and occurs once. Remove
successively \((3,1),(2,1),(1,1)\), then
\((-6,1),(-5,1),\ldots,(-1,1)\). Each removed ray is the sum of
its two then-adjacent rays. The inverse operation is the ordinary
two-coordinate blowup, with actual maps \((u,v)\mapsto(u,uv)\)
and \((u',v')\mapsto(u'v',v')\) in that retained lattice basis.
The result is precisely the three-ray fan
\((1,0),(0,1),(-1,0)\) of the native fixed sphere times its
parameter disc. Thus these are nine explicitly recorded interior
blowups. Their full normal monomial rows and original canonical
orders travel through the same displayed chart maps. This supplies
the resolved surface and normal input without substituting a
characteristic number for a chart or a clutch.



\par\bigskip\noindent\textbf{Preserved alternate construction: NC source-chart input}\par
This construction uses the eleven actual native-star charts NS.1-\/-6,

\par\bigskip\noindent\textbf{Preserved alternate construction: NC derivative chart count}\par
On each of the eleven actual resolved moving-star charts let

\par\bigskip\noindent\textbf{Preserved alternate construction: NC closed replacement source blowups}\par
actual nine-fold blowup of \(\mathbb P^1_Z\times \mathbb P^1_{\rm parameter}\) given by the
NS.23b fan and its recorded chronological coordinate blowup maps.
All nine centres lie in the original P disc.

\par\bigskip\noindent\textbf{Preserved alternate construction: NC unit/divisor scope}\par
among those units; their divisor orders stay in the eleven-chart
calculation.

\par\bigskip\noindent\textbf{Preserved alternate construction: original collision MN.1--2 and eleven-chart subtotal}\par
\par\bigskip\noindent\textbf{P3K.111 supporting proof 15: Two actual own-cap moving normal bundles}\par\nobreak
\par

\medskip\noindent\textbf{Selected-I3 \(M_{2}\): the two actual moving normal bundles and their own caps}

\par\nobreak

This calculation uses the executed eleven-chart collision star NS and
the original fourteen-factor reference RB. Each normal bundle is first
attached to its own held-Cox cap. The collision cap is the complete
NC.1--4 construction. The reference cap uses the same displayed ball
construction with the reference coefficients derived in MN.3 below.
The final section gives the exact dense comparison and its vertex
pole. The finite fan reset that compares the two whole completed
families is retained as a separate construction. No total return is
assigned here.

\par

\medskip\noindent\textbf{MN.1. Every collision normal coefficient and its divisor}

\par\nobreak

Keep NS.1--23, including the actual point modules, full
\(\eta,c_q,f_0\), the literal \(T_0\) alignment and the original
selected height. The moving surface has quotient tangent rays
\[
 (1,0),(3,1),(2,1),(1,1),(0,1),(-1,1),(-2,1),(-3,1),
                  (-4,1),(-5,1),(-6,1),(-1,0).
 \tag{MN.1}
\]
The two entries are \((\operatorname{ord}Z,\operatorname{ord}P)\).
The actual NS matrices give all adjacent cones, their inverse maps
and the normal columns. Let \(C_j\), \(-6\le j\le3\), denote
the compact curve belonging to \((j,1)\). Write \(F=\sum_jC_j\),
the actual divisor of P. The two remaining rays are the horizontal
zero and infinity sections of the fixed sphere.

In physical order \((-1,+1)\), let \(e_-,e_+\) be the coefficient
map from the executed normal coordinates to the regular native
coordinates. On the infinity charts the full units are
\((c_q/\eta,\eta)\); on the zero charts they are
\((c_qf_0/\eta,\eta/f_0)\). Their entire tangent monomials are
the rows NS.23 evaluated on the actual aligned cone. The derivative
has no suppressed source term: P and the native fixed-curve
coordinate have zero normal order, and p contains the product of
the two normals. Hence differentiating any of these full base units
has zero first normal term on the fixed surface.

The divisors \(D_- =\operatorname{div}(e_-)\) and
\(D_+ =\operatorname{div}(e_+)\), interpreted as the actual maps
from the executed lines to the pulled-back native lines, have the
following full coefficient list. The two horizontal orders are zero.

\begin{longtable}[]{@{}
  >{\raggedright\arraybackslash}p{(\columnwidth - 20\tabcolsep) * \real{0.1000}}
  >{\raggedleft\arraybackslash}p{(\columnwidth - 20\tabcolsep) * \real{0.0900}}
  >{\raggedleft\arraybackslash}p{(\columnwidth - 20\tabcolsep) * \real{0.0900}}
  >{\raggedleft\arraybackslash}p{(\columnwidth - 20\tabcolsep) * \real{0.0900}}
  >{\raggedleft\arraybackslash}p{(\columnwidth - 20\tabcolsep) * \real{0.0900}}
  >{\raggedleft\arraybackslash}p{(\columnwidth - 20\tabcolsep) * \real{0.0900}}
  >{\raggedleft\arraybackslash}p{(\columnwidth - 20\tabcolsep) * \real{0.0900}}
  >{\raggedleft\arraybackslash}p{(\columnwidth - 20\tabcolsep) * \real{0.0900}}
  >{\raggedleft\arraybackslash}p{(\columnwidth - 20\tabcolsep) * \real{0.0900}}
  >{\raggedleft\arraybackslash}p{(\columnwidth - 20\tabcolsep) * \real{0.0900}}
  >{\raggedleft\arraybackslash}p{(\columnwidth - 20\tabcolsep) * \real{0.0900}}@{}}
\toprule\noalign{}
\begin{minipage}[b]{\linewidth}\raggedright
Curve
\end{minipage} & \begin{minipage}[b]{\linewidth}\raggedleft
\(C_3\)
\end{minipage} & \begin{minipage}[b]{\linewidth}\raggedleft
\(C_2\)
\end{minipage} & \begin{minipage}[b]{\linewidth}\raggedleft
\(C_1\)
\end{minipage} & \begin{minipage}[b]{\linewidth}\raggedleft
\(C_0\)
\end{minipage} & \begin{minipage}[b]{\linewidth}\raggedleft
\(C_{-1}\)
\end{minipage} & \begin{minipage}[b]{\linewidth}\raggedleft
\(C_{-2}\)
\end{minipage} & \begin{minipage}[b]{\linewidth}\raggedleft
\(C_{-3}\)
\end{minipage} & \begin{minipage}[b]{\linewidth}\raggedleft
\(C_{-4}\)
\end{minipage} & \begin{minipage}[b]{\linewidth}\raggedleft
\(C_{-5}\)
\end{minipage} & \begin{minipage}[b]{\linewidth}\raggedleft
\(C_{-6}\)
\end{minipage} \\
\midrule\noalign{}
\endhead
\bottomrule\noalign{}
\endlastfoot
\(D_-\) & −4 & −4 & −4 & −4 & −4 & −5 & −6 & −5 & −4 & −3 \\
\(D_+\) & 5 & 5 & 5 & 5 & 5 & 5 & 4 & 3 & 2 & 1 \\
\end{longtable}

The companion recomputes each entry from the original four-cone
matrix and the full NS character, and verifies agreement from both
adjacent charts. This uses all eleven charts. On the common native
overlap the two line coefficients are carried by
\(a_\infty=Za_0\), \(c_\infty=Zc_0\); these are the exact native
line frames used to interpret the displayed divisors.

The original point/O identities are
\[
 V_{\rm pt}=-P^{-2j}E A_{\rm pt}
 =P^{-2j}c^{-1}E C_{\rm pt}
 =P^{-2j-2}(x-P^2)E_OC_O.
 \tag{MN.2}
\]
Thus their horizontal point or zero divisor is cancelled by its
actual Cartier generator before the normal derivative is taken.
The full products E and \(E_O\) and the extra \(P^{-2}\) at O
remain in the coefficient map. The map to the native coefficient
is exactly NS.16, not a degree-zero replacement of MN.2. On the
fixed locus the remaining units are the entire functions of
\((P,p=0)\) above. Each has its nonzero constant value and its
based logarithm on the P-disc. NC.4 gives the reciprocal extending
bundle homotopy in the held-Cox line projectors. The meromorphic
P and Z factors are excluded from that homotopy and have exactly
the orders in the table.

\par

\medskip\noindent\textbf{MN.2. Actual intersection matrix and the collision NC bundle}

\par\nobreak

The wall relation for the two endpoints of the compact chain is
\(v_{\rm left}+v_{\rm right}=v_j\), and at each of its other
vertices it is \(v_{\rm left}+v_{\rm right}=2v_j\).
Consequently the full ten-by-ten intersection matrix in the table\textquotesingle s
order is the tridiagonal matrix
\[
 M=\begin{pmatrix}
 -1&1&0&\cdots&0\\
 1&-2&1&\cdots&0\\
 0&1&-2&\ddots&\vdots\\
 \vdots&\vdots&\ddots&\ddots&1\\
 0&0&\cdots&1&-1
 \end{pmatrix}.
 \tag{MN.3}
\]
All adjacent determinants are one. This matrix has kernel generated
by the full fibre F. A generic fixed-sphere point divisor h meets
\(C_0\) once and every other \(C_j\) zero times, by the actual
projection given by the native Z character. Therefore
\[
 D_-^2=-5,\qquad D_-D_+=2,\qquad D_+^2=-4,
 \qquad hD_-=-4,\qquad hD_+=5.
 \tag{MN.4}
\]
These equations are direct products with MN.3 and its specified
marked section; no cap or degree is selected to obtain them.

The actual nine blowdowns in NS.23b remove, in order,
\[
 (3,1),(2,1),(1,1),(-6,1),(-5,1),(-4,1),(-3,1),(-2,1),(-1,1).
 \tag{MN.5}
\]
Each removed ray is the sum of its current neighbours. The two
coordinate blowup maps are the written NS maps \((u,v)\mapsto(u,uv)\)
and \((u',v')\mapsto(u'v',v')\). Let \(E_1,\ldots,E_9\) be
the pullbacks to the full surface of the exceptional curves contracted
at steps 1 through 9 of MN.5, respectively. Their matrix is \(-I_9\), and they are
orthogonal to the two pullback classes h and u from the fixed sphere
and the closed parameter sphere. Here h and u are classes on the explicit closed replacement
\(\overline F=F_{\rm star}-C\) constructed in NC.4, with
\(h^2=u^2=0\) and \(hu=1\); the original-disc divisor F
represents u there. NC.19--20 carries its rank-two bundle and tangent
stable bundle across the specified five-dimensional bordism to the
actual cap-completed surface \(F_{\rm star}+C_0\). All characteristic
numbers below are transferred along that bordism. Removing each actual ray and
subtracting the coefficients of its neighbours gives
\[
 \begin{aligned}
 D_-&=-4u+E_4+E_5+E_6-E_7-E_8,\\
 D_+&= 5u-E_4-E_5-E_6-E_7.
 \end{aligned}\tag{MN.6}
\]
This independently reproduces every product in MN.4 through the
actual nine coordinate blowups.

NC.17--20 constructs the cap product with the full resolved normal
map \(G_0\operatorname{diag}(\bar e_-,e_+)\), its other native
chart and the literal \(\widehat C B=N\) equality. Its explicit
five-dimensional bundle bordism replaces that cap by the opposite
fixed-sphere/parameter-disc face. The two positive line subspaces
have native classes \(+h,-h\). The coefficient map from the
executed negative line to its native line is \(e_-\), so its
unconjugated class is \(-h-D_-\); its complex conjugate has
class \(h+D_-\). The positive line has class \(-h-D_+\).
This fixes the signs from the coefficient direction. Equivalently,
NC.4\textquotesingle s actual opposite parameter-disc map assigns degree \(-n\)
to the target coefficient \(P^n\) on the boundary.

The two ordered line classes of the positive normal bundle on the
NC.4 closed replacement are thus
\[
 \begin{aligned}
 x&=h-4u+E_4+E_5+E_6-E_7-E_8,\\
 y&=-h-5u+E_4+E_5+E_6+E_7.
 \end{aligned}\tag{MN.7}
\]
Their exact products are
\[
 x^2=-13,\qquad xy=-3,\qquad y^2=6,
 \qquad c_1=x+y=-9u+2E_4+2E_5+2E_6-E_8,
\]
\[
                 c_1^2=-13,\qquad c_2=-3.
 \tag{MN.8}
\]
On the closed replacement the tangent Chern numbers are
\(c_1(T\overline F)^2=8-9=-1\) and
\(c_2(T\overline F)=4+9=13\), from the same nine coordinate
blowups of \(\mathbb P^1\times\mathbb P^1\). Hence
\(\int_{\overline F}p_1=-1-2(13)=-27\). The explicit oriented
bundle bordism NC.19--20 transfers \(c_1(E)^2,c_2(E)\) and
\(p_1(TF)\) to the specified actual NC-capped moving component;
its numbers are therefore exactly \((-13,-3,-27)\). The global orbifold
Stokes and return orientations are not assigned by MN.8.

\par



\par\bigskip\noindent\textbf{Preserved alternate construction: MN old collision-subtotal cross-reference}\par
The finite whole-pair continuation can use the exact q-preserving
Cox/ruling maps IF3.16--19 and the signed source/period resets
VI3.14--15, with the full current multipliers substituted. Those
maps, their native Hopf products, finite exterior and cap endpoints
must be composed to compare MN.8 with MN.14 in one prescribed
selected fan. This final composition is separate from the two
normal-bundle computations proved here.




\par\bigskip\noindent\textbf{Preserved alternate construction: MN retained actual-frame locator}\par
MN.2, and their remaining local factors are nonzero.

\par\bigskip\noindent\textbf{Preserved alternate construction: old IL.1 numerical declaration}\par
The twelve surfaces here are the actual \texttt{lower\_source\_boundary} face
orbits of the executed one-end atlas. They are closed against the
twelve scalar-source reference pieces in RB by the actual source and
line maps below. This gives the scoped physical characteristic pair
\[
 (k_M,\ell_M)=(-226,-1591)\quad\hbox{per lower double},\qquad
             (-2712,-19092)\quad\hbox{for all twelve}.
 \tag{IL.1}
\]
The moving selected surface, its Hopf cap, the whole reference fan
change RFB, and the final exterior/Jin transfer retain their own
specified arrows in the subsequent assembly.



\par\bigskip\noindent\textbf{Preserved alternate construction: IL normal row execution scope}\par
\((0,1)\) and \((1,0)\). These rows are retained for all 3,888 charts.

\par\bigskip\noindent\textbf{Preserved alternate construction: old IL.2--3 ordered array, 321 blowdowns and collision subtotal}\par
\medskip\noindent\textbf{IL.2. All source charts, normal walls and ordinary blowdowns}

\par\nobreak

Each of the twelve executed stars has 324 source charts and 325
ordered rays in \(0\leq x\leq3r\), with endpoints \((0,1),(3,1)\).
Every consecutive pair has determinant one and is an actual
executed source chart. The full endpoint lifts are
\[
       (0,12,1,0),\qquad(3,13,1,0)
                    =L(0,12,1,0).
 \tag{IL.9}
\]
The endpoint identification includes the coefficient \(p\) in
IL.3. It is evaluated at the actual target source coordinates.
The quotient identifies these ends and gives a central cycle of
324 compact curves. Its three surviving minimal source rays are
\((0,1),(1,1),(2,1)\); the last endpoint \((3,1)\) is their
periodic copy of \((0,1)\).

For a central ray \(r\), let \(r_-,r_+\) be its actual neighbours,
using the full inverse \(L\) at the first endpoint. The unimodular
four-dimensional matrix with columns
\((E_l,E_{l+1},r,r_-)\) gives the exact wall equation
\[
 -r_- -r_+=d_-E_l+d_+E_{l+1}+c\,r.
 \tag{IL.10}
\]
The last inverse-matrix entry is zero. The normal degrees on that
compact curve are \((d_-,d_+)\) and its source self-intersection is
\(c\). To see the direction, express the common meromorphic
coefficients IL.6 in the two smooth adjacent monomial charts.
Their coefficient transition is the corresponding integer power
of the inverse base coordinate. The divisor of that transition is
exactly \(d_-\), respectively \(d_+\), with the physical order in
IL.5. This also gives the reciprocal transition on the opposite
overlap. The complete nonzero unit factors are the restrictions of
IL.3; their two defining discs and reciprocal maps remain in that
calculation.

The executed blowdown sequence removes 321 rays from each source
star. At each step the removed primitive ray is the sum of its
two present neighbours, and its self-intersection is \(-1\).
The actual maps on its two affine charts are
\[
       (u,w)\longmapsto(u,uw),\qquad
       (u,w)\longmapsto(uw,w).
 \tag{IL.11}
\]
Their inverse maps on the complementary principal opens divide
the second, respectively first, coefficient by the retained
nonzero coordinate. The exceptional projective line is contracted
by these proper maps. The ordered list of the two neighbours,
the ray, both normal degrees, and both coefficient maps is saved
for every step, independently for every label \(l\).
Their composition is the actual source blowdown to the minimal
\(I_3\) source, with identity map on the retained outer boundary.
Every original source function is its pullback through the full
I2P/E81 map.

If \(d_\pm=N_\pm\cdot E\) at one step, tensoring \(N_\pm\) by
\(\mathcal O(d_\pm E)\) makes its exceptional degree zero. Its
two exact Cartier coefficient maps descend to the parent charts:
in the chart where \(E\) is \(u=0\), the map to the parent
coefficient is multiplication by \(u^{d_\pm}\), with reciprocal
module arrow on the common principal open. Consequently
\[
 N_\pm=\pi^*N_{\pm,0}\otimes\mathcal O(-d_\pm E),\qquad
 \Delta(c_-^2,c_-c_+,c_+^2)=(-d_-^2,-d_-d_+,-d_+^2).
 \tag{IL.12}
\]
Degrees on each adjacent strict curve add \(d_\pm\) in the
parent. The executable performs this integral operation together
with each source blowdown.

The final minimal-source degree vectors are
\[
 \deg N_-=(-1,0,1),\qquad \deg N_+=(-1,1,0)
                  \quad\hbox{on }(C_0,C_1,C_2).
 \tag{IL.13}
\]
The 321 chronological exceptional degree pairs consist of 49
copies of \((-1,0)\), 45 of \((0,-1)\), 133 of \((-1,-1)\),
and 94 of \((0,0)\). Summing the actual IL.12 maps gives
\[
       \Delta(c_-^2,c_-c_+,c_+^2)=(-182,-133,-178).
 \tag{IL.14}
\]
The source fan agrees for the twelve labels, while their full
normal lifts and period representatives remain separately stored.

\par

\medskip\noindent\textbf{IL.3. Relative lifts and the complete original tangent comparison}

\par\nobreak

With coefficient zero on \(C_0\), the integral source matrix is
\[
 I=\begin{pmatrix}-2&1\\1&-2\end{pmatrix}.
 \tag{IL.15}
\]
The degree vectors on \(C_1,C_2\) in IL.13 give the rational
relative lifts \(I^{-1}(0,1)^t=(-1/3,-2/3)^t\) and
\(I^{-1}(1,0)^t=(-2/3,-1/3)^t\). Their three pairings are
\[
        (-2/3,-1/3,-2/3).
 \tag{IL.16}
\]
These lifts retain the integral line bundles and their order-three
boundary characters. The full fibre \(F\), whose multiplicities
are the \(P\)-orders of the rays, is the intersection radical.
Both normal bundles have degree zero on \(F\). Adding its
coefficient to a relative lift therefore changes none of the
three pairings.

Combining IL.14 and IL.16 gives the collision-side pairings
\[
        (c_-^2,c_-c_+,c_+^2)=(-548/3,-400/3,-536/3).
 \tag{IL.17}
\]
An independent computation uses the complete 324-component cycle.
After the marked \(C_0\) is omitted its integral intersection
matrix is tridiagonal and negative definite. Exact successive
elimination of its diagonal pivots, followed by backward
substitution, solves the two full degree equations. Every row
is checked in the original matrix. The three resulting products
are exactly IL.17 for each of the twelve labelled stars.

The same maps IL.11 determine the tangent framing. Their
Jacobians have the single exceptional factor \(u\), respectively
\(w\), so \(K_{\widetilde S}=\pi^*K_S+E\), with
\(E^2=-1\) and zero pairing against the pullback class. The
exceptional projective line increases the relative Euler number
by one. Hence each step changes \(p_1=c_1^2-2c_2\) by \(-3\),
with the identity outer boundary map, and
\[
            \Delta\!\int p_1(TS)=-963.
 \tag{IL.18}
\]
The original volume in I2O is
\[
 -\frac{2J_{ts}}{a_T}P\,dP\wedge dv\wedge d\log X\wedge d\log Y.
 \tag{IL.19}
\]
Its entire unit, order \(2r+z-1\) on each ray, and all
coefficient derivatives are retained in the chart records.
IL.18 follows from the actual Jacobians relative to that fixed
outer volume and does not change its boundary marking.

\par



\par\bigskip\noindent\textbf{Preserved alternate construction: old IL.39--40 closed numerical pairings and tangent subtotal}\par
Gluing through
IL.30, using the cube connections and the retained fibre
gauges, subtracts IL.37 from IL.17 and gives
\[
       (c_-^2,c_-c_+,c_+^2)_{\rm double}
                           =(-182,-134,-178).
 \tag{IL.39}
\]
For the tangent bundle, the minimal collision source is
carried to the entire reference body by the explicit
relative source flow. Its derivative and inverse determine
matching boundary tangent connections. Their identical
opposite-collar interpolation cancels. The remaining
contribution is exactly the 321 actual blowups in IL.18:
\[
                 \int_{\rm double}p_1(TS)=-963.
 \tag{IL.40}
\]


\par\bigskip\noindent\textbf{Preserved alternate construction: IL residue orientation reference}\par
the complex fixed-surface orientation already used in
IL.10--40.

\par\bigskip\noindent\textbf{Preserved alternate construction: old IL.42 lower totals and absent execution provenance}\par
Substitution of the three complete normal pairings and
the tangent value gives
\[
 k_M=-182+134-178=-226,\qquad
 \ell_M=-963-182-268-178=-1591.
 \tag{IL.42}
\]
All twelve fixed-face orbits have the same pair by their
separately executed full maps, hence IL.1. Their distinct
labels and period representatives are retained.

The primary \path{code/part3/e6_selected_i3_m2_lower_normals.py}
records all 3,888 inverse chart matrices, original
source/normal/volume rows, all central wall maps, and
3,852 chronological ordinary blowdowns. It also solves
both full intersection systems independently. The initial
collision-only execution passed 66,128 exact assertions;
its complete source/output is preserved under
\path{records/attempts/e6_selected_i3_m2_lower_collision_only/.}
The companion relative certificate checks the full
IL.26--38 period and cube-frame identities, the exact
reference subtraction, the physical Euler expansion,
and each of the twelve separate resulting rows.



\par\bigskip\noindent\textbf{Preserved alternate construction: IL final reference scope}\par
by \((-72,0)\). IL.1 uses the actual RB \(T_1\)
reference pieces and its own cap maps.

\par\bigskip\noindent\textbf{Preserved alternate construction: ASI2 old common and endpoint execution input}\par
Take precisely the executed common I2F fan and its proper zero,
one and infinity lifts I2ZC/I2L. The complete common boundary has
136,442 primitive regular three-cones. The actual parent semigroup
maps, full character inverses, paired period facets and protected
source cuts are the maps in CFL.2--3. Every refined cone is the
specified proper toric subdivision of its raw parent; the original
monomial relations pull back identically. 

\par\bigskip\noindent\textbf{Preserved alternate construction: ASI2 old fixed inventory and moving collision roots}\par
The fixed-face inventory examines every regular
cone and every actual paired period facet. Its
minimal nonzero-weight faces, exact stabilizer
words and complete stars give the following
disjoint list:

\begin{itemize}

\item
  106,166 compact generic curves, with all three
  complete endpoint normal and determinant maps;
\item
  15,138 compact generic surfaces with their
  saturated quotient stars and 52 actual blowups;
\item
  zero corner interiors: 6,163 points, 72 curves
  and 17 complete compact surfaces;
\item
  one corner interiors: 4,764 points and 2,340 curves;
\item
  infinity corner interiors: 7,304 points and 701 curves;
\item
  the twelve lower/reference source doubles of IL;
\item
  the single NS moving face and the RB moving
  reference face, with their respective caps.
\end{itemize}

The first five rows are exactly CFL.1--9. All have
their complete period alignments and full analytic
normal units. Every generic component has already
been completed across all three ends; the listed
corner-interior tags are disjoint from those
generic boundary pieces. The twelve lower
stabilizers are their actual source \(L\) groups.
The moving face is the single orbit identified
by the full \(T_0=hB^{-1}\) map NS.5. No additional
fixed orbit remains in the exhaustive face
inventory or in the exterior cylinder.

Their first five rows give \((2723,19124)\).
IL gives the next row \((-2712,-19092)\).
For the moving row, retain the actual positive
normal bundle \(E=\overline{N_-}\oplus N_+\).
NC.19--20 is an explicit five-dimensional bundle
bordism from its own cap-completed surface to
the opposite-disc closed replacement. On that
replacement the nine actual source blowups
give the MN.7 roots
\[
 x=h-4u+E_4+E_5+E_6-E_7-E_8,\quad
 y=-h-5u+E_4+E_5+E_6+E_7,
 \tag{ASI2.40}
\]
with \(h^2=u^2=0,hu=1,E_iE_j=-\delta_{ij}\).
Thus \(c_1(E)^2=-13,c_2(E)=-3,p_1(TF)=-27\),
transferred through that actual bordism.
The reference roots are \(h-12u,-h-12u\);
its own cap gives \(c_1^2=c_2=p_1(TF)=0\).


\par\bigskip\noindent\textbf{Preserved alternate construction: ASI2 old moving pair and old inner ledger}\par
Hence the moving collision-minus-
own-reference cap pair is \((-10,-28)\).
The full inner ledger with the RB \(T_1\) reference is
\[
 (2723,19124)+(-2712,-19092)+(-10,-28)=(1,4).
 \tag{ASI2.42}
\]
This ledger includes its entire reference exactly once.



\par\bigskip\noindent\textbf{Preserved alternate construction: ASI2 finite slice scope}\par
finite covering factor in CFL and IL.

\par\bigskip\noindent\textbf{Preserved alternate construction: ASI2 orientation scope}\par
fixed-surface orientations used in CFL, IL and
MN.

\par\bigskip\noindent\textbf{Preserved alternate construction: ASI2.7 exact inherited-map clarification}\par
The outer fibres are the original
oriented spheres of F41/F46 with the
proved \(M_{2}\) good cubic cap. An explicit
path retains the selected colliding
branch and \(r_{\rm src}\) real from
\(0\) to \(1/4-s_*\), with
\(0<s_*<1/12\). K96 gives the full
\(M_{2}\) source/Cartier/good-root family at
the initial unselected III value.
The remaining interval avoids the
source collisions and selected cusp.
The exact SI2 Miller, point/O, source
and cubic-root maps continue along
that interval with the prescribed
finite character and integral
annular reference. A finite compact
parameter cover gives a common strict
contraction bound. Its invariant
horizontal fields give the actual
equivariant fibre diffeomorphism,
with backward inverse. The original
Hopf tube and primitive attaching
maps of F41/F35 travel through this
same source/cap family; their
Gysin and van Kampen maps retain
the sphere topology for the stated
selected fan.
